% Appendices A to D of the New In ML 2026 camera-ready paper.

\section{Proofs and positioning for \texorpdfstring{Section~\ref{sec:nonexp}}{Section 3}}\label{app:early}

\paragraph{Integral minimisers suffice.} On bipartite systems, $\ell_0$ and $\ell_1$ distances to
the optimal face may be computed over its integral points (Section~\ref{sec:prelim}, footnote), and
the small-instance figures of this paper are computed that way. The face may contain non-integral
points, but none is ever strictly closer. Let $y$ be integral, let $f$ minimise the Hamming distance
to $y$ over the optimal face $F$, and let $S$ be the coordinates where they agree. Substituting
$y_S$ into the covering system leaves a system on the remaining coordinates. Its matrix is a
submatrix of a bipartite incidence matrix, hence totally unimodular, and its right-hand side is
integral. So $\{z : z \ge 0,\ z_i + z_j \ge w_{ij},\ z_S = y_S\}$ is an integral polyhedron, and
so is its optimal face for the integral objective $\sum_i z_i$, which is nonempty (it contains $f$)
and has value $\opt$. Any integral point of that face agrees with $y$ on all of $S$, so it is an
integral minimiser. For $\ell_1$ the substitution is not needed. Fix a sign pattern
$I \subseteq V$. On the region $\{z : z_i \ge y_i\ (i \in I),\ z_i \le y_i\ (i \notin I)\}$ the
objective $\|y - z\|_1$ is linear, and the region intersected with the covering system is again a
totally unimodular system with integral right-hand side. Its optimal face for $\sum_i z_i$ is $F$
intersected with the region, and it is integral. So every region that meets $F$ contains an
integral $\ell_1$-minimiser, and hence so does $F$. For $\ell_\infty$ and $\ell_\infty^\pm$ this
reduction does not apply, and there we minimise over the full face directly, fractional points
included.

\begin{proof}[Proof of Corollary~\ref{cor:lip}]
Clamping is componentwise $1$-Lipschitz, so assume $x, x' \ge 0$, and run the two scans in
lockstep. At a constraint $(i,j)$ there are three cases. If neither iterate fires, nothing changes.
If both fire, each maps $(z_i, z_j) \mapsto (w_{ij} - z_j,\, w_{ij} - z_i)$, so on those two
coordinates the difference of the two iterates is swapped and negated. That is a signed
permutation, hence an isometry of every $\ell_p$. If exactly one fires, say the iterate $z$ while
$z'$ does not, then $z'_i + z'_j \ge w_{ij}$. This is the only property of $y$ used in the proof of
Theorem~\ref{thm:nonexp}, so that proof applies verbatim with $y := z'$. Composing over the scan
gives the claim.
\end{proof}

Theorem~\ref{thm:nonexp} is the case $x' = y$ feasible, where $\repair_\sigma(y) = y$. We state it
separately because the runtime bounds compose with this form.

In Corollary~\ref{cor:dinitz}, Algorithm~1 of \citeauthor{dinitz2021} is their repair (a one-sided
walk), not \citeauthor{chen2022}'s solver of the same name.

\begin{proof}[Proof of Corollary~\ref{cor:dinitz}]
The inequality inside the proof of their Lemma~12 bounds the iteration count by $g(\hat{y}')$ for
any feasible integer dual $\hat{y}'$, in particular for the one their Algorithm~1 produces. (Their
Lemma~12 then weakens $g(\hat{y}')$ to $\|\hat{y}' - y^*\|_1$; we skip that step.) Their Lemma~6
bounds that algorithm's total movement by twice the minimum feasible correction, and
$\delta = (\hat{y}-y^*)^+$ is feasible for the correction problem: since $y^*$ is feasible,
$\hat{y}_i + \hat{y}_j - c_{ij} \le (\hat{y}_i - y^*_i) + (\hat{y}_j - y^*_j) \le
\delta_i + \delta_j$. Hence the movement is at most $2\|(\hat{y}-y^*)^+\|_1$, so
$\sum_i \hat{y}'_i \ge \sum_i \hat{y}_i - 2\|(\hat{y}-y^*)^+\|_1$ and
\[
g(\hat{y}') = \sum_i y^*_i - \sum_i \hat{y}'_i
\;\le\; \|(y^*-\hat{y})^+\|_1 + \|(\hat{y}-y^*)^+\|_1 = \|\hat{y} - y^*\|_1 .
\]
Alternatively, replacing their repair by the packing mirror of Definition~\ref{def:greedy} gives the
same bound directly from Theorem~\ref{thm:pack}.
\end{proof}

\begin{remark}
\citeauthor{dinitz2021}'s Theorem~7 bounds the post-repair $\ell_1$ distance and
Corollary~\ref{cor:gap} the post-repair gap, so the two meet only in the end-to-end bound, which is
where Corollary~\ref{cor:dinitz} compares them.
\end{remark}

\paragraph{Relation to known results.}
(i) Relaxation and reflection methods give the $\ell_2$ case of Theorem~\ref{thm:nonexp} (Section~\ref{sec:nonexp}).
(ii) Order-preserving operators and (iii) tropical projection do not apply. (iv) Isotonic
regression is the closest precedent.

\emph{(ii) Order-preserving operators.} \citet{crandall1980} show that on a domain closed under
componentwise maximum a sum-preserving operator is order preserving if and only if it is
$\ell_1$-non-expansive. This does not apply: $\repair_\sigma$ is not sum-preserving (it strictly
increases $\sum_i y_i$ on every infeasible input), and it is not order preserving.

\emph{(iii) Tropical projection.} The residuated projector $P(y) = \max\{v \in V : v \le y\}$ onto a
max-plus semimodule~\citep{cohen2004} is a classical sup-norm
non-expansion~\citep[Prop.~2]{crandall1980}. The $\ell_1$ analogue fails: on the max-plus line
$V = \{(t,t,t)\} \subset \R^3$, $P(y) = (\min_i y_i)\mathbf{1}$ has $\ell_\infty$ Lipschitz
constant $1$ but $\ell_1$ constant $3 = n$ (take $y = (0,1,1)$, $z = (1,1,1)$). In our setting no
exact residuated projector exists, since the packing set is not max-closed, and its natural
analogue (Proposition~\ref{prop:fail}(a)) is not non-expansive either.

\emph{(iv) Isotonic regression.} \citet{yang2019} prove that the $\ell_2$ isotonic projection onto
the monotone cone $\{x_1 \le \dots \le x_n\}$ is non-expansive in every permutation-invariant norm,
by averaging a point with its image under an adjacent transposition. The same averaging reaches
$\dfeas$. Its constraint normals are $e_i + e_j$ and
$e_i$, so on each piece of the exact $\ell_2$ projection onto $\dfeas$ the linear part is an
average of signed permutation matrices, and that projection is $1$-Lipschitz in every $\ell_p$ as
well. The metric property therefore does not single out $\repair_\sigma$; what distinguishes it is
that it has the property while being local, integral, one-pass and linear-time, and that its step,
the full residual, is the largest common step that keeps it (Appendix~\ref{app:delicate}).

\paragraph{A name collision.} \citet{kumabe2025} study ``Lipschitz continuous algorithms for
covering problems'', including vertex cover. Their Lipschitz constant perturbs the input weights
and measures earth-mover distance between output distributions. In that framework no deterministic
algorithm is Lipschitz for most problems they consider, and their lower bound is $\Omega(|V|)$ for
exact vertex-cover algorithms. Our constant perturbs the starting point and measures how far a
single deterministic output moves, so their statements and ours are unrelated.

\paragraph{The twisted lattice.} Proposition~\ref{prop:lattice} rules out one possible explanation
of the operator map (Table~\ref{tab:zoo}, Appendix~\ref{app:delicate}): that the covering direction
has no canonical repair because $\dfeas$ is not a lattice. It is a lattice under the twisted order,
so the discrete-convex-analysis machinery of \citet{sakaue2022,sakaue2023} formally applies to it.
The twist works only in the bipartite case: the flip $z_b = -y_b$ is a proper $2$-colouring of the
constraint graph. For a non-bipartite covering system (weighted vertex cover on a graph with an odd
cycle) no sign flip of this kind exists, and the $L^\natural$-convex route is not available.
Theorem~\ref{thm:nonexp} does not need bipartiteness.

\begin{proposition}[The feasible set is a twisted lattice]\label{prop:lattice}
$\dfeas$ is not closed under the componentwise minimum. It \emph{is} closed under the twisted
operations $(\max$ on $A$, $\min$ on $B)$ and $(\min$ on $A$, $\max$ on $B)$; equivalently, the
isometry $z_b := -y_b$ turns $\dfeas$ into a box intersected with a difference-constraint system,
which is an $L^\natural$-convex (alcoved) polyhedron.
\end{proposition}
\begin{proof}
For the failure, a single constraint with $w = 10$ and the points $(10,0)$, $(0,10)$: their minimum
$(0,0)$ violates $y_1 + y_2 \ge 10$. For the twisted closure, take the operation $(\max$ on $A$, $\min$ on $B)$ and any constraint
$(a,b)$. Without loss of generality $\min(v_b, v'_b) = v_b$, and then
$\max(u_a,u'_a) + \min(v_b,v'_b) \ge u_a + v_b \ge w_{ab}$. The operation
$(\min$ on $A$, $\max$ on $B)$ is symmetric. Under $z_b = -y_b$ these are the two lattice
operations of the difference-constraint system $y_a - z_b \ge w_{ab}$.
\end{proof}

\section{Non-expansion is fragile: proofs and the full operator map}\label{app:delicate}

\begin{proof}[Proof of Proposition~\ref{prop:fail}]
For (a), a small explicit witness on a $2\times2$ packing instance is
$c = \left(\begin{smallmatrix}1&3\\5&0\end{smallmatrix}\right)$,
$\hat{y} = (-1, 6 \mid -1, 3)$, $y = (-3, 2 \mid -3, -3)$. The repair drives the second left coordinate
from $6$ down to $-3$, past $y$'s value of $2$, and the $\ell_1$ distance to $y$ rises from $14$ to
$15$.

For the unbounded constant, take all capacities zero and the feasible point $y = 0$. For (a), take
a $k \times 2$ instance with $\hat y = (\mathbf{0} \mid 1, 0)$. The $A$-pass lowers every $u_a$ to
$\min_b(0 - v_b) = -1$, and the $B$-pass then leaves $v$ alone, so the output is
$(-\mathbf{1} \mid 1, 0)$ at $\ell_1$ distance $k+1$ from $y$, against $\|\hat y - y\|_1 = 1$. For (b),
take the $k \times 1$ analogue $\hat y = (\mathbf{0} \mid 1)$. Every constraint has excess $1$, and the
integral split puts all of it on the $A$-side ($\lceil 1/2 \rceil = 1$,
$\lfloor 1/2 \rfloor = 0$), so each $u_a$ falls to $-1$ while $v$ never moves: again distance $k+1$
against $1$. Both outputs are feasible, and $|V| = k+2$ and $k+1$ respectively.
\end{proof}

Operator (a) is the order-theoretic way to project down to feasibility, the analogue of the
residuated projector of tropical algebra. In the packing direction, however, maximal feasible
points below $\hat{y}$ need not be unique, because the feasible set is not closed under
componentwise maximum (single constraint $y_1 + y_2 \le 0$, $\hat{y} = (1,1)$: both $(1,-1)$ and
$(-1,1)$ are maximal). So every projection of this kind implicitly fixes an order of sides. (The
Euclidean projection onto the same set is non-expansive in every $\ell_p$, since its normals are
$-(e_a + e_b)$, but it is a global quadratic program with non-integral output; see
Appendix~\ref{app:early}, item (iv).)

Operator (b) splits the excess as evenly as integrality allows. It fails too, because that split is
uneven. In the covering notation of Theorem~\ref{thm:nonexp}, raising both endpoints by $r/2$ is a
non-expansion in every $\ell_p$: the worst case $s_j = r - s_i$ contributes
$2\,|s_i - r/2|^p - |s_i|^p - |s_i - r|^p \le 0$ by midpoint convexity of $|t|^p$. But that step
is not integral. A common integral step keeps both properties. Any common step $h$
with $r/2 \le h \le r$ (for integral data this includes $\lceil r/2 \rceil$) is non-expansive in
every $\ell_p$ and satisfies the constraint it fires. The full residual is the largest such step.
Since $2h \ge r$, the step restores the fired constraint, and since coordinates only increase, one
pass suffices. With $h$ in place of $r$ in $f$, which is still non-increasing, the worst case
$s_j = r - s_i$ contributes $|s_i - h|^p + |s_i - r + h|^p - |s_i|^p - |s_i - r|^p \le 0$ by
convexity, because the two new points lie between the two old ones and have the same sum. A step
$h > r$ expands: on $y_1 + y_2 \ge r$ from $0$, it increases the distance to
$\tfrac{r}{2}(e_1 + e_2)$.

\paragraph{Local and global operators.} By Proposition~\ref{prop:lattice} (Appendix~\ref{app:early}),
lattice structure does not explain the operator map. The relevant distinction
is between \emph{local} and \emph{global} operators. \citeauthor{sakaue2022} define their
repair as an $\ell_\infty^{\pm}$-projection, that is, as an $\arg\min$, which need not be unique. As
representative they exhibit a \emph{uniform shift} of all $n$ coordinates ($+\varepsilon/2$ on one
side, $-\varepsilon/2$ on the other, in their orientation $s_i - t_j \ge w_{ij}$). That
representative has $\ell_1$ displacement $\Theta(n\varepsilon)$ and is not $\ell_1$-non-expansive.
But the $\arg\min$ also contains sparse minimisers, so $\ell_\infty^\pm$-optimality does not force
the large $\ell_1$ displacement. The repair of Definition~\ref{def:greedy}
is local: it touches only the endpoints of violated constraints, and Corollary~\ref{cor:lip} shows
that it is $1$-Lipschitz in every $\ell_p$ at once.

We transplant the $\ell_\infty^\pm$-projection to the covering orientation.
(i) \citet[Lemma~2]{sakaue2022} exhibit the uniform shift as one $\ell_\infty^\pm$-projection for
the difference orientation. (ii) Their rounding is to the nearest integer with $0.5$ rounded down,
which preserves feasibility in their orientation but not in ours: for an integral clamped
prediction it fails at every odd $\varepsilon$, because both endpoints of a maximally violated
constraint round down and give back the last unit of slack. We round up instead (our choice). (iii) Their runtime bounds
are parameterised by an $\ell_\infty$-type distance from the prediction
($O(T_{\mathrm{loc}}\|p^*-\hat p\|_\infty)$ in \citet[Thm.~1(2)]{sakaue2022},
$\widetilde O(T_{\mathrm{loc}}\,\bar\mu(\hat p;g))$ in \citet[Thm.~1(ii)]{sakaue2023}), and their
solver is a steepest descent whose iteration count is bounded by that same distance. Both papers
also cap the worst case by running a standard algorithm in parallel. (iv) The two-sided norm is
evaluated in covering coordinates; in the difference coordinates of \citet{sakaue2022} the face
point is itself $\ell_\infty^\pm$-optimal on the family of Observation~\ref{obs:shift}, so the
forcing claim below is stated for the plain sup norm $\|\cdot\|_\infty$.
Observation~\ref{obs:shift} shows a property of the norm: a repair can be optimal
in $\ell_\infty$ and still wrong in every coordinate, so an $\ell_\infty$ guarantee is of no use
to a bound stated in $\ell_0$.

\begin{lemma}[Agreement with the optimal face is limited by $\opt$]\label{lem:budget}
Let $z \in \R^V_{\ge 0}$ and put $s = \min_{i \in V} z_i$. If $s > 0$, then the $\ell_0$
distance from $z$ to the full optimal face is at least $|V| - \lfloor \opt / s \rfloor$.
\end{lemma}
\begin{proof}
Fix an optimal $y$ and let $S = \{ i : z_i = y_i \}$. Every optimal $y$ has $y \ge 0$ and
$\sum_{i} y_i = \opt$, so $\sum_{i \in S} z_i = \sum_{i \in S} y_i \le \opt$. Each summand on the
left is at least $s$, so $s\,|S| \le \opt$.
\end{proof}

\begin{observation}[An $\ell_\infty$-optimal repair can be wrong in every coordinate]\label{obs:shift}
Transplant the $\ell_\infty^\pm$-projection to the covering orientation: if the clamped
prediction's largest violation is $\varepsilon > 0$, raise every coordinate by $\varepsilon/2$ and
round up. The output $z$ has $z_i \ge \lceil \varepsilon/2 \rceil \ge 1$ everywhere, so by
Lemma~\ref{lem:budget} its $\ell_0$ distance to the optimal face is at least
$(|A|+|B|) - \lfloor \opt / \lceil \varepsilon/2 \rceil \rfloor$. On the following family the
output is wrong in every coordinate. For $d \ge 2$ and $W \ge 2$ take $|A| = 1$, $|B| = d$ and
$w_{ab} = W$ for every $b$. The optimal face is the single point $(W \mid 0, \dots, 0)$; the prediction
$\hat y = 0$ lies at $\ell_0$ distance $1$ from it; and the repair outputs
$\lceil W/2 \rceil \mathbf{1}$, which differs from it in all $|A| + |B|$ coordinates. Taking $k$ disjoint such stars keeps the face a single point, keeps the post-repair
face distance at $|A| + |B|$, and leaves the prediction's own at $|A|$. On this family the
$\ell_0$ error is forced by the norm, whatever the representative:
\emph{every} feasible integral $z$ attaining $\min \|z - \hat y\|_\infty$ (equal here to the
covering-coordinate $\|z - \hat y\|_\infty^\pm$, since $z \ge \hat y$) differs from the
optimal face in all $|A| + |B|$ coordinates.
\end{observation}
\begin{proof}
For one star, let $y$ be feasible and $t = y_a$. If $t \ge W$ then $\sum_i y_i \ge t \ge W$, with
equality only at $t = W$ and $y \equiv 0$ on $B$. If $t < W$ then $y_b \ge W - t > 0$ for all $d$
coordinates of $B$, so $\sum_i y_i \ge t + d(W-t) = dW - (d-1)t > W$ using $t < W$ and $d \ge 2$.
Hence $\opt = W$, attained only at $(W \mid 0,\dots,0)$. The prediction $\hat y = 0$ agrees with it on
all of $B$, so its face distance is $1$; the largest violation is $\varepsilon = W$; the output
is $\lceil W/2 \rceil \mathbf{1}$, and $\lceil W/2 \rceil \notin \{0, W\}$ because $W \ge 2$. For
$k$ disjoint stars the cross constraints have weight $0$ and are implied by $y \ge 0$, so the
instance decouples and the face is the product of the per-star faces.

For the last claim, $z \ge 0 = \hat y$ gives $\|z - \hat y\|_\infty = \max_i z_i$, and
feasibility with integrality gives $\max(z_a, z_b) \ge \lceil W/2 \rceil$ for each $b$; the uniform
shift attains $\lceil W/2 \rceil$, so that is the minimum. Any $z$ attaining it has
$z_a \le \lceil W/2 \rceil < W$ and therefore $z_b \ge W - \lceil W/2 \rceil = \lfloor W/2 \rfloor
\ge 1$ for every $b$, both strict inequalities using $W \ge 2$. So $z$ agrees with
$(W \mid 0,\dots,0)$ nowhere.
\end{proof}

The transplant keeps the $\ell_\infty$-optimality. A constraint violated by $\varepsilon$
forces some endpoint to move by at least $\varepsilon/2$ in any feasible repair, hence by at least
$\lceil\varepsilon/2\rceil$ in any integral repair, the kind the downstream bounds are stated for.
The uniform shift realises displacement $\lceil\varepsilon/2\rceil$, so it is optimal in the norm
its guarantee is stated in.

Lemma~\ref{lem:budget}, however, is vacuous once $\opt \gg |V|\lceil\varepsilon/2\rceil$, which is the
regime of the random instances of Section~\ref{sec:exp} and Appendix~\ref{app:experiments}. The
family is also an extreme case: on those random instances, $\ell_\infty$-optimality does not force
an error in every coordinate. The \emph{endpoint shift} raises only the endpoints of violated constraints by
$\varepsilon/2$ and rounds up. It is also $\ell_\infty$-optimal, with the same $\ell_\infty$
displacement as the uniform shift on all $120$ draws, yet it incurs only $0.34$ to $0.76$ of the
uniform shift's post-repair $\ell_0$ error in every cell. Section~\ref{sec:exp} shows that the
uniform shift is also costly for the solver: warm-starting through it can take several times as many
dual adjustments as the row-reduction start.

\paragraph{The full operator map.} Table~\ref{tab:zoo} collects what is known about every repair in
this paper.

\makeatletter\setlength{\@fptop}{0pt}\makeatother
\begin{table}[!htbp]
\centering
\caption{What each repair guarantees in $\ell_1$ and in $\ell_0$; ``$?$'' marks an open cell. The
walk and uniform-shift rows are our covering transplants of operators of \citet{dinitz2021} and
\citet{sakaue2022,sakaue2023}; the rounding up is ours, and the entries of these rows are our analysis.
Lower-bound and tightness entries are proved on covering instances. In the packing setting we claim
only the factor-$2$ upper bound of Theorem~\ref{thm:composed} and the balanced lower bound of
Theorem~\ref{thm:balancedlb}.}
\label{tab:zoo}
\footnotesize
\setlength{\tabcolsep}{4pt}\begin{tabular}{@{}p{2.8cm} p{3.8cm} p{6.6cm}@{}}
\toprule
Repair & $\ell_1$ & $\ell_0$ face distance\\
\midrule
\raggedright Greedy two-sided, full residual & \raggedright non-expansive in every $\ell_p$ & \raggedright no additive constant under any order (Thm.~\ref{thm:noadditiveclass}); on bipartite covering instances the class of all $\ell_1$-non-expansive repairs has constant exactly $3/2$ (Thm.~\ref{thm:threehalves})\tabularnewline
\raggedright \quad forward scan & \raggedright non-expansive & \raggedright $\Omega(d)$ (Thm.~\ref{thm:order})\tabularnewline
\raggedright \quad sorted (weakly-descending) & \raggedright non-expansive & \raggedright $\le \|y^*{-}\hat y\|_0 + 2\min(|A|,|B|){-}1$ (Thm.~\ref{thm:sharpwarm}; attained at every $\min(|A|,|B|)$, Prop.~\ref{prop:sharpattained}; ${+}1$ at $|A|{=}1$, Prop.~\ref{prop:sortedwarm}); and $\ge \|y^*{-}\hat y\|_0 + \min(|A|,|B|){-}1$ on an explicit balanced family (Thm.~\ref{thm:balancedlb}): no additive constant\tabularnewline
\raggedright \quad cover-directed (Thm.~\ref{thm:coverscan}) & \raggedright non-expansive & \raggedright $\le \|y^*{-}\hat y\|_0 + \tau(G(p)) \le 2\|y^*{-}\hat y\|_0$, tight in the limit for some cover and order (Prop.~\ref{prop:chainfam}); not additive (Thm.~\ref{thm:noadditiveclass})\tabularnewline
\raggedright $\ell_1$-nearest face point (Prop.~\ref{prop:faceadditive}; consults the optimal face) & \raggedright bound toward optimal duals only: $\le 2\|\hat y{-}y^*\|_1$ to every optimal $y^*$, meeting \citeauthor{chen2022}'s composition hypothesis; \emph{not} known non-expansive toward all of $\dfeas$ & \raggedright additive with constant $0$: the output lies on the face\tabularnewline
\raggedright Greedy, integral split & \raggedright \emph{not} non-expansive; no finite constant (Prop.~\ref{prop:fail}) & \raggedright $?$\tabularnewline
\raggedright Alternating projection (A-first) & \raggedright \emph{not} non-expansive; no finite constant (Prop.~\ref{prop:fail}) & \raggedright $?$\tabularnewline
\raggedright One-sided walk & \raggedright \emph{not} non-expansive (ratio $2$ witness, Sec.~\ref{sec:exp}); movement $\le 2\|\hat y{-}y^*\|_1$, distance $\le 3\|\hat y{-}y^*\|_1$ & \raggedright $\le \|y^*{-}\hat y\|_0 + \tau(G(p)) \le 2\|y^*{-}\hat y\|_0$ under every tie-break, no additive constant (Prop.~\ref{prop:walkbound}); the factor~$2$ tight in the limit (Prop.~\ref{prop:ladderfam}); can lose to the sorted scan under fixed tie-breaks (Sec.~\ref{sec:exp}); lowest among the repairs in Tables~\ref{tab:exp} (reference-optimum $\ell_0$) and~\ref{tab:solver} (dual adjustments)\tabularnewline
\raggedright Uniform shift ($\ell_\infty^\pm$-projection, rounded up) & \raggedright $\ell_\infty^\pm$-optimal in covering coordinates; $\ell_1$ displacement $\Theta(n\varepsilon)$ & \raggedright face-$\ell_0 \ge (|A|{+}|B|) - \lfloor\opt/\lceil\varepsilon/2\rceil\rfloor$; wrong in every coordinate on an explicit family (Obs.~\ref{obs:shift})\tabularnewline
\bottomrule
\end{tabular}
\end{table}

\clearpage
\section{Order-dependence: the family in full}\label{app:order}

This appendix proves Theorem~\ref{thm:order}, computes the quantities on which its two orders agree
and disagree, and tabulates the family for $d \le 6$ (Table~\ref{tab:order}). Throughout, the instance has $|A| = 1$ with
left coordinate $u$ and right coordinates $v_0, \dots, v_{d-1}$. The forward scan processes the
constraints in index order, which on this family is lightest first, and the reverse scan heaviest
first.

\begin{proof}[Proof of Theorem~\ref{thm:order}]
Take $w_{0k} = k+1$ for $k = 0,\dots,d-1$ and $\hat{y} = 0$. Writing $\mathrm{val}(t) = t +
\sum_k (k+1-t)^+$ for the best dual with $u = t$, we have $\mathrm{val}(t) = t + (d - t) = d$ for
$t \in [d-1, d]$ and $\mathrm{val}(t) > d$ for $t < d-1$, so $\opt = d$ and the integral optimal
face consists of exactly the two points $y^{(d)} = (d \mid \mathbf{0})$ and
$y^{(d-1)} = (d-1 \mid e_{d-1})$. (The optimum is not unique, consistent with the face formula in the
proof of Proposition~\ref{prop:sortedwarm}, where $W_2 = d-1$.) The prediction $\hat{y} = 0$
differs from $y^{(d)}$ in one coordinate.

Forward scan: at step $k$ the residual is $(k+1) - k - 0 = 1$, so $u$ rises by $1$ and $v_k$ becomes
$1$. The output is $u = d$, $v \equiv 1$. It differs from $y^{(d)}$ in the $d$ right coordinates,
and from $y^{(d-1)}$ in $u$ and the $d-1$ right coordinates below the top, so its face distance is
$d$.

Reverse scan: the first step $k = d-1$ has residual $d$, giving $u = d$ and $v_{d-1} = d$; all
later residuals are non-positive. The output differs from $y^{(d)}$ in one coordinate.

Both outputs have $\sum_i y_i = 2d$, hence the gap $d$.
\end{proof}

\paragraph{The two $\ell_1$ distances.} The two $\ell_1$ equalities also follow by direct computation. Both outputs lie
at $\ell_1$ displacement $2d$ from the prediction (forward: $d$ unit raises of $u$ and $d$ right
coordinates raised to $1$; reverse: $u$ and $v_{d-1}$ each raised by $d$). On the continuous face
$\{(t \mid (d-t)e_{d-1}) : t \in [d-1,d]\}$, writing $s = d - t \in [0,1]$, the forward output is at
$\ell_1$ distance $s + (d-1) + (1-s) = d$ and the reverse output at $s + (d-s) = d$, for every $t$.
The prediction $\hat y = 0$ is at $t + s = d$. In either order, the repair leaves the $\ell_1$
distance to the optimal face unchanged. The gap and both $\ell_1$ distances therefore cannot tell
the orders apart, and neither can the realised $\ell_\infty$ displacement from the prediction, which is
$d$ for both.

\paragraph{The $\ell_\infty$ reversal.} The $\ell_\infty$ distance to the face does tell the orders
apart, but it ranks them the wrong way. With $s = d - t$ as above, the forward output
$(d \mid \mathbf{1})$ is at $\max(s,1,|1-s|) = 1$ and the reverse output $(d \mid d\,e_{d-1})$ at
$\min_t \max(s,t) = d-1$. The same two values hold in the $\ell_\infty^\pm$ norm of
\citet{sakaue2022} evaluated in covering coordinates; in their difference coordinates the forward
output is still at $1$ and the reverse output is at $d$, so the reversal persists. For $d \ge 3$
the $\ell_\infty$ ranking is therefore strict and reverses the $\ell_0$ ranking. An iteration bound
parameterised by the $\ell_\infty^\pm$ distance from the starting point to the optimal set, the form
of \citet[Prop.~2]{sakaue2022}, would prefer the order that an $\ell_0$-parameterised bound must
avoid.

\begin{table}[ht]
\centering
\caption{The family of Theorem~\ref{thm:order}. The two orders agree on the dual gap, the realised $\ell_1$ displacement and the
$\ell_1$ face distance, and the last equals the prediction's own.}
\label{tab:order}
\small
\setlength{\tabcolsep}{4.5pt}
\begin{tabular}{rrrrrrrrrrr}
\toprule
& \multicolumn{3}{c}{$\ell_0$ to face} & \multicolumn{2}{c}{$\ell_\infty$ to face}
& \multicolumn{2}{c}{dual gap} & \multicolumn{2}{c}{$\ell_1$ to face} & realised\\
\cmidrule(lr){2-4}\cmidrule(lr){5-6}\cmidrule(lr){7-8}\cmidrule(lr){9-10}
$d$ & pre & fwd & rev & fwd & rev & fwd & rev & fwd & rev & $\ell_1$ disp.\\
\midrule
2 & 1 & 2 & 1 & 1 & 1 & 2 & 2 & 2 & 2 & 4\\
3 & 1 & 3 & 1 & 1 & 2 & 3 & 3 & 3 & 3 & 6\\
4 & 1 & 4 & 1 & 1 & 3 & 4 & 4 & 4 & 4 & 8\\
5 & 1 & 5 & 1 & 1 & 4 & 5 & 5 & 5 & 5 & 10\\
6 & 1 & 6 & 1 & 1 & 5 & 6 & 6 & 6 & 6 & 12\\
\bottomrule
\end{tabular}
\end{table}

\paragraph{The $\ell_1$ dependence is not slack.} On this family the prediction is at face
distance $1$, so in the chain of Section~\ref{sec:composition} a bound in the prediction's $\ell_0$
error would give $O(m\sqrt{n} + m + n\log n)$. The forward scan yields post-repair $\ell_0 = d$,
which is the dense regime. (This compares the parameter only: with $|A| = 1 < |B|$ the family has no
perfect matching, so it is not an instance of \citet[Thm.~3.2]{chen2022}.) Both orders are allowed
by Theorem~\ref{thm:nonexp} and agree on the gap and both $\ell_1$ distances. We do not prove that no
$\ell_1$-based analysis could select between them: that would need a lower bound over a defined
class of analyses, which we leave open.

\section{The sorted scan in full}\label{app:sorted}

This appendix proves the results of Section~\ref{sec:sorted} on the sorted repair
$\repair_\downarrow$ (Definition~\ref{def:sorted}) and gives the cold-start results in full:
$2\min(|A|,|B|)$ is not an envelope (Proposition~\ref{prop:coldrefuted}), and the exact envelope is
$\lfloor\tfrac52\min(|A|,|B|)\rfloor - 1$ for $\min(|A|,|B|) \ge 3$
(Theorem~\ref{thm:coldfivehalves}).

\subsection{Cold start}

\begin{proof}[Proof of Proposition~\ref{prop:sortedcold}]
Induction over the firings. Suppose the statement holds for all earlier ones and consider a firing at
$(a,b)$. If $a$ were an endpoint of an earlier firing $(a,b_0)$, the inductive hypothesis gives
$y_a = w_{ab_0}$, and descending order gives $w_{ab_0} \ge w_{ab}$, so the residual at $(a,b)$ is
$w_{ab} - y_a - y_b \le w_{ab} - w_{ab_0} \le 0$ and no firing occurs, a contradiction. The same
argument applies to $b$. Hence both endpoints are untouched and therefore $0$, the residual is
$w_{ab}$, and both coordinates become $w_{ab}$. So the firings have pairwise distinct endpoints,
that is, they form a matching, and a matching in $K_{A,B}$ has at most $\min(|A|,|B|)$ edges.
\end{proof}

An exhaustive check covers all $2{,}360{,}750$ weight matrices with entries in $\{0,\dots,4\}$ at the
shapes $1{\times}3$, $1{\times}4$, $2{\times}2$, $2{\times}3$, $2{\times}4$ and $3{\times}3$, under
every weakly-descending tie-break order ($134{,}081$ cold outputs across all tie-breaks). The cold
face distance there is at most $2\min(|A|,|B|) - 1$, attained at each of the six shapes. The bound
does not extend: a $3\times4$ instance (weights $\le 6$; $\le 3$ under a fixed tie-break) ends at
face distance $6 = 2\min(|A|,|B|)$ under every weakly-descending order.
Proposition~\ref{prop:sortedcold} counts moved coordinates and does not bound the face distance,
since the output can also disagree with the face on unmoved coordinates. The natural link between
the two, an optimal dual supported on the endpoints of the fired matching, fails in general.

\begin{proposition}[The cold start can miss the face on every coordinate]\label{prop:coldrefuted}
\begin{enumerate}
\item[(a)] For every instance, every weakly-descending order and every optimal $f$, the cold output $q$
satisfies $\|q - f\|_0 \le |\mathrm{supp}\, f| + |M_f| \le 3\min(|A|,|B|)$, where $M_f$ is the
fired matching; and trivially $\|q - y^*\|_0 \le |A| + |B|$.
\item[(b)] On
$w = \left(\begin{smallmatrix} 12&11&0&7&10\\ 11&2&8&11&5\\ 10&11&8&12&10\\ 9&0&7&5&1 \end{smallmatrix}\right)$
the optimal face is the single point $f = (10,9,10,7 \mid 2,1,0,2,0)$, every weakly-descending
order fires the same matching $\{(1,1),(2,3),(3,4),(4,5)\}$ and outputs
$q = (12,8,12,1 \mid 12,0,8,12,1)$, and $q$ disagrees with $f$ on all $9 = |A| + |B|$ coordinates:
$\|q - y^*\|_0 = 9 > 8 = 2\min(|A|,|B|)$ under every tie-break.
\item[(c)] The cold face distance exceeds $2\min(|A|,|B|)$ by a constant factor.
The $k$-fold direct sum of (b) is a $4k \times 5k$ instance with cold face distance
$9k = |A| + |B|$ against $2\min(|A|,|B|) = 8k$, a ratio of $9/4$ to $\min(|A|,|B|)$. A
$6 \times 8$ instance has cold face distance $14 = |A| + |B|$ under the lexicographic tie-break, so
its $k$-fold direct sum gives $\|q - y^*\|_0 / \min(|A|,|B|) = 7/3$ at arbitrarily large
$\min(|A|,|B|)$; a $5 \times 8$ instance with a unique optimal dual and cold face distance $11$
under the lexicographic tie-break (weights $\le 7$) gives $11/5$.
Theorem~\ref{thm:coldfivehalves} gives the exact envelope
$\lfloor \tfrac{5}{2}\min(|A|,|B|) \rfloor - 1$ for $\min(|A|,|B|) \ge 3$, which the
$4 \times 5$, $5 \times 8$ and $6 \times 8$ witnesses attain at their own $\min(|A|,|B|)$.
\end{enumerate}
\end{proposition}
\begin{proof}
(a) Every fired vertex $v$ has $q_v = w_{ab} > 0$ for its fired pair and every unfired vertex has
$q_v = 0$, so $q$ and $f$ disagree on an unfired vertex only if $f_v > 0$, and on a fired pair
$(a, b)$ on at most two coordinates, at most one of which can have $f = 0$ since
$f_a + f_b \ge w_{ab} > 0$. Hence $\|q - f\|_0 \le |\mathrm{supp}\, f \setminus V(M_f)| +
|\mathrm{supp}\, f \cap V(M_f)| + |M_f| = |\mathrm{supp}\, f| + |M_f|$; and $\mathrm{supp}\, f
\subseteq V(M^*)$ for every maximum-weight matching $M^*$ by complementary slackness, so the two
terms are at most $2\min(|A|,|B|)$ and $\min(|A|,|B|)$ respectively. (b) is a finite check:
$f$ is feasible and $\sum f = 41 = w(M^*)$ for $M^* = \{(1,1),(2,4),(3,2),(4,3)\}$. The
uniqueness of $f$ and the single cold output over all $2{,}304$ weakly-descending orders of the
positive weights are verified exactly by computer, together with (c) (Appendix~\ref{app:repro}). The face,
the cold scan and the Hamming distance all decompose over the zero-separated blocks of a direct
sum, which gives (c).
\end{proof}

In (b) the greedy matching is not maximum-weight ($33 < 41$), it fires all four rows, and every
fired column receives a weight that no optimal dual agrees with, so no coordinate is left in agreement. Below
four rows the envelope $2\min(|A|,|B|)$ holds. At four rows a $4 \times 5$ witness with weights
$\le 6$ already exceeds it under the lexicographic tie-break. The upper bound of (a) is not tight: a finer count lowers it
to $\lfloor \tfrac52 \min(|A|,|B|) \rfloor$ for every optimal dual, and a shift within the face
lowers the face distance by one more, to a bound that the families of
Proposition~\ref{prop:coldfamilies} attain.

\begin{theorem}[The cold envelope is exactly $\lfloor \tfrac52\min(|A|,|B|) \rfloor - 1$ for $\min(|A|,|B|) \ge 3$]\label{thm:coldfivehalves}
Let $m = \min(|A|,|B|)$, fix any weakly-descending order, and let $q$ be the output of
$\repair_\downarrow$ from $\hat y = 0$. Write $F$ for the optimal face.
\begin{enumerate}
\item[(i)] For every optimal $f$, $\|q-f\|_0 \le \lfloor \tfrac52 m \rfloor$. For a fixed $f$ this
  is sharp: for every $m \ge 1$ some instance has the same cold output under every
  weakly-descending order and an optimal $f$ with $\|q-f\|_0 = \lfloor \tfrac52 m \rfloor$.
\item[(ii)] If $m \ge 2$, the face distance satisfies
  \[
    \|q-y^*\|_0 \;=\; \min_{f \in F} \|q-f\|_0 \;\le\; \Bigl\lfloor \tfrac{5}{2} m \Bigr\rfloor - 1 ,
  \]
  that is, $\tfrac52 m - 1$ for even $m$ and $\tfrac{5m-3}{2}$ for odd $m$.
\item[(iii)] Let $E(m)$ be the largest cold face distance $\|q-y^*\|_0$ over all instances with
  $\min(|A|,|B|) = m$ and all weakly-descending orders. Then $E(1) = 1$, $E(2) = 3$, and
  \[
    E(m) \;=\; \Bigl\lfloor \tfrac{5}{2} m \Bigr\rfloor - 1 \qquad \text{for every } m \ge 3,
  \]
  attained by instances whose cold output is the same under every weakly-descending order. In
  particular the bound of (ii) is attained for every $m \ge 3$, and not at $m = 2$.
\end{enumerate}
\end{theorem}

\begin{proof}
Write $m=\min(|A|,|B|)$ and assume without loss of generality $|A| \le |B|$ (the
scan and the claim are symmetric in the two sides).  Let the fired pairs be
$(a_i,c_i)$ with weights $g_i = w_{a_i c_i} > 0$, $i = 1,\dots,k$; they form a
matching (Proposition~\ref{prop:sortedcold}), $q$ equals $g_i$ on both endpoints
of the $i$-th pair and $0$ elsewhere, and $k \le m$.  Fix an optimal $f$ and a
maximum-weight matching $M^*$, $|M^*| =: m' \le m$.  By complementary slackness
every edge of $M^*$ is tight at $f$ ($f_a + f_b = w_{ab}$) and $M^*$ saturates
$\mathrm{supp}\, f$.

\emph{Classification.}  Since $f_{a_i}+f_{c_i}\ge g_i>0$, at most one endpoint of
a fired pair carries $f=0$.  Call the pair \emph{BP} if both endpoints are
positive, \emph{ZC} if $f_{c_i}=0$ (then $f_{a_i}\ge g_i$; \emph{ZC$^>$} when the
inequality is strict, \emph{ZC$^=$} otherwise), and \emph{ZA} if $f_{a_i}=0$
(then $f_{c_i}\ge g_i$; \emph{ZA$^>$}, \emph{ZA$^=$} likewise), with counts
$\mathrm{bp},\mathrm{zc}^>,\mathrm{zc}^=,\mathrm{za}^>,\mathrm{za}^=$, so
$k=\mathrm{bp}+\mathrm{zc}^>+\mathrm{zc}^=+\mathrm{za}^>+\mathrm{za}^=$.  Let
$u_r$ (resp.\ $u_c$) be the number of unfired rows (columns) with $f>0$; since
all fired rows are distinct, $u_r \le |A|-k = m-k$.

\emph{Disagreement count.}  On a fired pair, $q$ equals $g_i$ twice; a ZC$^=$
pair agrees at its row ($f_{a_i}=g_i$) and a ZA$^=$ pair at its column, so BP,
ZC$^>$, ZA$^>$ pairs contribute at most $2$ disagreements and ZC$^=$, ZA$^=$
pairs at most $1$.  Every unfired vertex has $q=0$, so unfired vertices
contribute exactly $u_r+u_c$.  Hence
\begin{equation}\label{eq:dcount}
  \|q-f\|_0 \;\le\; 2\,\mathrm{bp} + 2\,\mathrm{zc}^> + 2\,\mathrm{za}^> +
  \mathrm{zc}^= + \mathrm{za}^= + u_r + u_c .
\end{equation}

\emph{Firing lemma.}  No edge $e$ of positive weight can have both of the
following properties: each endpoint of $e$ is either unfired or fired at a pair
of weight strictly less than $w_e$; and $e$ is not itself a fired pair.  Indeed,
in a weakly-descending order every fired pair of weight $< w_e$ is processed
strictly after $e$, and a vertex moves only when its pair fires; so when $e$ is
processed both endpoints still stand at $0$, the residual is $w_e > 0$, and $e$
fires.

\emph{Hosting.}  Consider an unfired column $d$ with $f_d>0$.  It is
$M^*$-saturated; let $(r,d) \in M^*$, a tight edge of weight $f_r + f_d > 0$.
If $r$ is unfired, the firing lemma applies to $(r,d)$ directly: impossible.
If $r=a_i$ with pair $i$ of type ZC, then $w_{a_i d} = f_{a_i} + f_d \ge g_i +
f_d > g_i$, so the pair of $a_i$ is lighter than the edge while $d$ is unfired:
the firing lemma again.  Hence every unfired positive column is
$M^*$-matched to the row of a BP or ZA pair.  Likewise, a ZA$^>$ column $c_i$
is $M^*$-matched to some row $r \neq a_i$ (matching to $a_i$ would force
$f_{c_i} = w_{a_i c_i} = g_i$).  Its tight edge has weight $f_r + f_{c_i} >
f_{c_i} > g_i$, heavier than the pair of $c_i$; if $r$ is unfired, or if
$r = a_j$ of type ZC (whence $w \ge g_j + f_{c_i} > g_j$), the firing
lemma forbids the edge.  So every ZA$^>$ column is also $M^*$-matched to
a BP or ZA row.  Since $M^*$ is a matching and there are only $\mathrm{bp} +
\mathrm{za}^> + \mathrm{za}^=$ BP-or-ZA rows,
\begin{equation}\label{eq:host}
  u_c + \mathrm{za}^> \;\le\; \mathrm{bp} + \mathrm{za}^> + \mathrm{za}^=,
  \qquad\text{i.e.}\qquad
  u_c \;\le\; \mathrm{bp} + \mathrm{za}^= .
\end{equation}

\emph{Two counts.}  First, substituting \eqref{eq:host} and $u_r\le m-k$ into
\eqref{eq:dcount} and rounding $\mathrm{zc}^=$, $\mathrm{za}^=$ up to twice
themselves:
\begin{equation}\label{eq:countone}
  \|q-f\|_0 \;\le\; 3\,\mathrm{bp} + 2(\mathrm{zc}^>+\mathrm{zc}^=+\mathrm{za}^>+\mathrm{za}^=) + (m-k)
           \;=\; m + k + \mathrm{bp} .
\end{equation}
Second, the rows of $\mathrm{supp}\, f$ (there are
$\mathrm{bp}+\mathrm{zc}^>+\mathrm{zc}^=+u_r$ of them) and the columns of
$\mathrm{supp}\, f$ ($\mathrm{bp}+\mathrm{za}^>+\mathrm{za}^=+u_c$) are
saturated by the matching $M^*$, so each count is at most $m' \le m$; adding
these to $\mathrm{zc}^>+\mathrm{za}^> \le k - \mathrm{bp}$ and comparing with
\eqref{eq:dcount}:
\begin{equation}\label{eq:counttwo}
  \|q-f\|_0 \;\le\; (\mathrm{bp}+\mathrm{zc}^>+\mathrm{zc}^=+u_r) +
  (\mathrm{bp}+\mathrm{za}^>+\mathrm{za}^=+u_c) + (\mathrm{zc}^>+\mathrm{za}^>)
           \;\le\; 2m + k - \mathrm{bp} .
\end{equation}
Adding \eqref{eq:countone} and \eqref{eq:counttwo}:
$2\|q-f\|_0 \le 3m + 2k \le 5m$, and $\|q-f\|_0$ is an integer. That is the bound in (i).

\emph{Sharpness of (i).} On
$w = \left(\begin{smallmatrix} 5&0&4\\ 4&1&0 \end{smallmatrix}\right)$ the pair of weight $5$ fires
first, both pairs of weight $4$ meet one of its endpoints, and the pair of weight $1$ then fires, so
every weakly-descending order outputs $q = (5,1 \mid 5,1,0)$. The dual $f = (3,2 \mid 2,0,1)$ is
feasible with $\sum f = 8 = \opt$ (the matching $\{(1,3),(2,1)\}$), and it differs from $q$ in all
$5 = \lfloor \tfrac52 \cdot 2 \rfloor$ coordinates. On $w = (2\ \ 1)$ every order outputs
$q = (2 \mid 2, 0)$, and the optimal $f = (1 \mid 1, 0)$ differs from it in
$2 = \lfloor \tfrac52 \rfloor$ coordinates. The optimal face, the cold scan and the Hamming distance
all decompose over zero-separated direct sums (as in the proof of
Proposition~\ref{prop:coldrefuted}(c)), so $\lfloor m/2 \rfloor$ copies of the first block, together
with one copy of the second when $m$ is odd, give for every $m \ge 1$ an order-free instance and an
optimal $f$ with $\|q-f\|_0 = \lfloor \tfrac52 m \rfloor$. Hence no count that holds for every
optimal $f$ can improve (i). The improvement in (ii) comes from choosing $f$ in the face, through
the shifts \eqref{eq:shiftone} and \eqref{eq:shifttwo} below.

\emph{Slack identities.} For an optimal $f$ write
$\rho_f = \mathrm{bp}+\mathrm{zc}^>+\mathrm{zc}^=+u_r$ and
$\kappa_f = \mathrm{bp}+\mathrm{za}^>+\mathrm{za}^=+u_c$ for the numbers of rows and of columns of
$\mathrm{supp}\, f$, and $\Delta_f \ge 0$ for the amount by which \eqref{eq:dcount} is strict.
Substituting $k = \mathrm{bp}+\mathrm{zc}^>+\mathrm{zc}^=+\mathrm{za}^>+\mathrm{za}^=$ into the
right-hand side of \eqref{eq:dcount} gives, with no inequality dropped,
\begin{align}
  \|q-f\|_0 &= m+k+\mathrm{bp} - s_1(f), \notag\\
  s_1(f) &= \Delta_f + (m-k-u_r) + (\mathrm{bp}+\mathrm{za}^=-u_c) + \mathrm{zc}^= ,
  \label{eq:sone}\\[2pt]
  \|q-f\|_0 &= 2m+k-\mathrm{bp} - s_2(f), \notag\\
  s_2(f) &= \Delta_f + (m-\rho_f) + (m-\kappa_f) + \mathrm{zc}^= + \mathrm{za}^= .
  \label{eq:stwo}
\end{align}
For \eqref{eq:stwo}, note that the first expression in \eqref{eq:counttwo},
$\rho_f + \kappa_f + \mathrm{zc}^> + \mathrm{za}^>$, is not a bound but a rewriting of the
right-hand side of \eqref{eq:dcount}. Every bracket is non-negative, by $u_r \le m-k$,
\eqref{eq:host} and $\rho_f, \kappa_f \le m' \le m$. So \eqref{eq:sone} and \eqref{eq:stwo} are
\eqref{eq:countone} and \eqref{eq:counttwo} with their slacks written out. (In
\eqref{eq:countone} only $\mathrm{zc}^=$ is rounded up, since \eqref{eq:host} has already
doubled $\mathrm{za}^=$.) Adding the two, $s_1(f)+s_2(f) = 3m+2k-2\|q-f\|_0$.

\emph{Proof of (ii).} Let $m \ge 2$, with $|A| = m \le |B|$ as above. Call an optimal $f$ a
\emph{minimiser} if it minimises $\|q-f\|_0$ over the optimal face. Minimisers may be taken
integral (Section~\ref{sec:prelim}, justified in Appendix~\ref{app:early}), and every move below
changes $f$ by an integer amount, so integrality is kept; the argument also works verbatim over the
full face. Let $d_0 = \|q-y^*\|_0$ be the common value, and suppose, for a contradiction, that
$2d_0 \ge 5m-1$.

\emph{Step 1: every row fires.} Since $5m-1 \le 2d_0 \le 3m+2k$, we have
$k \ge m - \tfrac12$, so $k = m = |A|$. Every row is fired, so $u_r = 0$, the bracket $m-k-u_r$
in \eqref{eq:sone} vanishes, and $\rho_f = \mathrm{bp}+\mathrm{zc}^>+\mathrm{zc}^=
= m-\mathrm{za}^>-\mathrm{za}^=$. The number $k$ depends on the order and not on $f$, so every
minimiser $f$ satisfies
\[
  s_1(f)+s_2(f) = 5m-2d_0 \le 1, \qquad
  s_2(f) = \Delta_f + \mathrm{za}^> + 2\,\mathrm{za}^= + \mathrm{zc}^= + (m-\kappa_f).
\]

\emph{Step 2: where a minimiser agrees with $q$.} Let $f$ be a minimiser; one of $s_1(f)$,
$s_2(f)$ is $0$. If $s_2(f) = 0$, then $\Delta_f = \mathrm{zc}^= = \mathrm{za}^= = 0$. If
$s_1(f) = 0$, then $\Delta_f = \mathrm{zc}^= = 0$ by \eqref{eq:sone}, and $s_2(f) \le 1$ forces
$\mathrm{za}^= = 0$, because $\mathrm{za}^=$ enters $s_2$ with coefficient $2$. (Count one alone
does not give this: from $s_1 = 0$ it gives only $u_c = \mathrm{bp}+\mathrm{za}^=$.) So in every
case $\Delta_f = 0$ and $\mathrm{zc}^= = \mathrm{za}^= = 0$. Then \eqref{eq:dcount} is an
equality and every fired pair is BP, ZC$^>$ or ZA$^>$, so it disagrees at both endpoints. There
is no unfired row, and an unfired column disagrees exactly when $f$ is positive there. Hence a
minimiser $f$ agrees with $q$ exactly on
\[
  Z_f = \{ d \in B : d \text{ unfired},\ f_d = 0 \}, \qquad d_0 = |A|+|B|-|Z_f| .
\]

\emph{Step 3: the unfired-column lemma.} If $a_i$ is a fired row and $d$ an unfired column, then
$w_{a_i d} \le g_i$. Otherwise $a_i$ is fired at a pair lighter than $(a_i,d)$, $d$ is unfired,
and $(a_i,d)$ is not a fired pair, which the firing lemma forbids.

\emph{Step 4: Shift I, a minimiser with at most $m-1$ positive columns.} Let $f$ be a minimiser
with $\kappa_f = m$, and let $P$ be its set of positive columns, so $|P| = m = |A|$. For
$0 \le e \le e^* := \min_{b \in P} f_b$ put
\begin{equation}\label{eq:shiftone}
  f^{(e)} = f + e\,(\mathbf{1}_A - \mathbf{1}_P) .
\end{equation}
The sum is unchanged because $|A| = |P|$. A pair $(a,b)$ with $b \in P$ keeps $f_a+f_b$, and a
pair with $b \notin P$ gains $e$. So $f^{(e)}$ is feasible, and it is non-negative since
$e \le e^*$. Hence $f^{(e)}$ is optimal. The columns of $Z_f$ are zero, so they lie outside $P$
and do not move. Thus $f^{(e)}$ still agrees with $q$ on $Z_f$, so $\|q-f^{(e)}\|_0 \le d_0$,
with equality by minimality. Hence $f^{(e^*)}$ is a minimiser, and the column where $e^*$ is
attained has dropped to $0$, so $\kappa_{f^{(e^*)}} \le m-1$. Fix a minimiser $f$ with
$\kappa_f \le m-1$.

\emph{Step 5: the shape it is forced into.} By Step 1, $s_2(f) \ge m-\kappa_f \ge 1$, so
$s_2(f) = 1$ and $s_1(f) = 0$. In particular $2d_0 = 5m-1$ and $m$ is odd, which already rules
out even $m$. For odd $m$ ($m \ge 3$), $s_2(f) = 1$ forces $\kappa_f = m-1$ and
$\mathrm{za}^> = \mathrm{za}^= = 0$, so $\rho_f = m$ and every fired pair is BP or ZC$^>$. And
$s_1(f) = 0$ forces $u_c = \mathrm{bp}+\mathrm{za}^= = \mathrm{bp}$. Hence
$m-1 = \kappa_f = \mathrm{bp}+u_c = 2\,\mathrm{bp}$, so $\mathrm{bp} = \tfrac{m-1}{2} \ge 1$.

\emph{Step 6: Shift II, and the contradiction.} Let $X$ be the set of BP rows and $Y$ the set of
unfired columns with $f > 0$, so $|X| = \mathrm{bp} = u_c = |Y| \ge 1$. For
$0 \le e \le \min_{b \in Y} f_b$ put
\begin{equation}\label{eq:shifttwo}
  f^{(e)} = f + e\,(\mathbf{1}_X - \mathbf{1}_Y) .
\end{equation}
The sum is unchanged. A pair $(a,b)$ with $b \notin Y$ loses nothing, and a pair with $a \in X$,
$b \in Y$ keeps $f_a+f_b$. Now let $b \in Y$ and $a \notin X$. Then $a$ is fired (Step 1) and not
BP, so its pair $(a,c_i)$ is ZC$^>$ and $f_a > g_i \ge w_{ab}$ by Step 3; hence
$f^{(e)}_a + f^{(e)}_b \ge f_a \ge w_{ab}$. Non-negativity holds for $e \le \min_Y f$, so
$f^{(e)}$ is optimal. It still agrees with $q$ on $Z_f$, which meets neither $X$ nor $Y$. At
$e = \min_Y f$ some $y \in Y$ reaches $0 = q_y$, which is a new agreement. So
$\|q-f^{(e)}\|_0 \le d_0-1$, contradicting minimality.

Therefore $2d_0 \le 5m-2$. That is (ii): $d_0 \le \tfrac52 m - 1$ for even $m$, and, since $d_0$
is an integer, $d_0 \le \tfrac{5m-3}{2}$ for odd $m$.

\emph{Proof of (iii), small $m$.} First let $m = 1$, say $A = \{a\}$. The first positive pair
processed has both endpoints at $0$ and fires, so $g_1 = \max w = \opt$. Then
$(\opt \mid 0,\dots,0)$ is optimal and differs from $q$ only at $c_1$, so $E(1) \le 1$; on
$w = (2\ \ 1)$ no optimal dual equals $q$, since $\sum q = 4 > 2 = \opt$. (So the bound of (ii)
also holds at $m = 1$.)

Now let $m = 2$, say $A = \{a_1,a_2\}$ with first fired pair $(a_1,c_1)$, so again $g_1 = \max w$.
We give an optimal $f$ that agrees with $q$ at $a_2$ and at every column other than $c_1$ and $c_2$,
so $\|q-y^*\|_0 \le 3$. If $k = 0$ all weights vanish and $f = q = 0$.

If $k = 1$, then $a_2$ is unfired, so by the firing lemma the only pair at $a_2$ that can be
positive is $(a_2,c_1)$. Let $\beta = w_{a_2c_1}$, and let $f$ be $\opt-\beta$ at $a_1$, $\beta$
at $c_1$ and $0$ elsewhere. The pair $(a_1,c_1)$ is covered because $\opt \ge g_1$. A pair
$(a_1,b)$ with $b \ne c_1$ is covered because $\{(a_1,b),(a_2,c_1)\}$ is a matching, and
$(a_2,c_1)$ is covered with equality. The sum is $\opt$, and $f \ge 0$ since
$\opt \ge \beta$. This $f$ differs from $q$ at most at $a_1$ and $c_1$.

If $k = 2$, let $(a_2,c_2)$ be the second fired pair, $\beta = w_{a_2c_1}$ and
$\beta' = \max(\beta, g_2)$. Put $f_{a_1} = \opt-\beta'$, $f_{a_2} = g_2$,
$f_{c_1} = \beta'-g_2$, and $0$ on every other column; the sum is $\opt$. The pair $(a_1,c_1)$
has $f_{a_1}+f_{c_1} = \opt-g_2 \ge g_1$, since the fired pairs form a matching. For $(a_1,b)$
with $b \ne c_1$, $w_{a_1b}+\beta \le \opt$ because $\{(a_1,b),(a_2,c_1)\}$ is a matching.
Also $w_{a_1b}+g_2 \le \opt$: for $b \ne c_2$ because $\{(a_1,b),(a_2,c_2)\}$ is a matching, and
for $b = c_2$ because $w_{a_1c_2} \le g_1$. Hence $f_{a_1} \ge w_{a_1b}$. Next,
$(a_2,c_1)$ has $f_{a_2}+f_{c_1} = \beta' \ge \beta$, $(a_2,c_2)$ has $g_2$, and
$(a_2,d)$ with $d$ unfired has $g_2 \ge w_{a_2d}$ by the unfired-column lemma (Step 3).
Non-negativity holds because $\opt \ge \max(\beta, g_2)$. So $E(2) \le 3$.

The block $w = \left(\begin{smallmatrix} 5&0&4\\ 4&1&0 \end{smallmatrix}\right)$ above attains
$3$. Its maximum matching $\{(1,3),(2,1)\}$ is unique, so every optimal $y$ is tight on it and
vanishes on column $2$. Hence $y_{a_1} \le 4 < 5$, $y_{c_1} \le 4 < 5$ and
$y_{c_2} = 0 < 1$ against $q = (5,1 \mid 5,1,0)$. The optimal dual $(4,1 \mid 3,0,0)$ attains
$3$. So $E(2) = 3$, below the value $4$ that (ii) allows.

\emph{Proof of (iii), attainment for $m \ge 3$.} For $m \ge 4$ this is
Proposition~\ref{prop:coldfamilies}. For $m = 3$ take the $3\times4$ witness of this section,
$w = \left(\begin{smallmatrix} 4&6&5&3\\ 5&4&3&0\\ 4&4&0&1 \end{smallmatrix}\right)$. The pairs
$(1,2)$, $(2,1)$, $(3,4)$ of weights $6, 5, 1$ form a matching, and every other positive pair
meets one of them at a strictly heavier pair. By the dominance argument in the proof of
Proposition~\ref{prop:coldfamilies}, every weakly-descending order outputs
$q = (6,5,1 \mid 5,6,0,1)$. The maximum matching is unique:
$M^* = \{(1,3),(2,1),(3,2)\}$, of weight $14$. So every optimal $y$ has $y_{c_4} = 0$,
$y_{a_1} = 5 - y_{c_3} \le 5$, $y_{c_2} = 4 - y_{a_3} \le 4$ and $y_{a_2}+y_{c_1} = 5$. That
forces four disagreements: at $c_4$ ($0 \ne 1$), at $a_1$ ($\le 5 < 6$), at $c_2$
($\le 4 < 6$), and at least one of $a_2, c_1$, since both would need the value $5$. So $y$ agrees
with $q$ on at most one of $a_2, c_1$, and possibly at $a_3$ and $c_3$. The two cases below show
that an agreement at $c_3$ or at $a_3$ excludes every other agreement, so $y$ agrees with $q$
on at most one coordinate.
\begin{itemize}
\item If $y_{c_3} = 0$, then $y_{a_1} = 5$, and $(1,2)$ gives $y_{c_2} \ge 1$, so
  $y_{a_3} \le 3$. Also $(2,3)$ gives $y_{a_2} \ge 3$, so $y_{c_1} \le 2$, and then $(3,1)$
  gives $y_{a_3} \ge 2$. So $a_3$ and $c_1$ disagree. And $a_2$ agrees only if $y_{a_2} = 5$,
  that is $y_{c_1} = 0$; then $(3,1)$ forces $y_{a_3} \ge 4$, contradicting $y_{a_3} \le 3$.
\item If $y_{a_3} = 1$, then $y_{c_2} = 3$, and $(3,1)$ gives $y_{c_1} \ge 3$. Then $(2,2)$ gives
  $y_{a_2} \ge 1$, so $y_{c_1} \le 4 < 5$. And $a_2$ agrees only if $y_{a_2} = 5$, that is
  $y_{c_1} = 0 < 3$. So $a_2$ and $c_1$ both disagree. And $c_3$ disagrees too, because
  $y_{c_3} = 0$ would force $y_{a_3} \ge 2$ by the previous case.
\end{itemize}
So $\|q-y^*\|_0 \ge 7 - 1 = 6 = \lfloor \tfrac{15}{2} \rfloor - 1$, with equality by (ii).
\end{proof}

\begin{proposition}[Extremal families]\label{prop:coldfamilies}
For $p \ge 2$ let $E_p$ be the instance with rows $R_1,\dots,R_p,S_1,\dots,S_p$ and columns
$H_1,\dots,H_p,T_1,\dots,T_p,U_1,\dots,U_p$ (so $m = 2p$ and $|B| = 3p$). Set
\[
  r_i = 2p+i,\quad h_i = p+i,\quad s_i = 2i,\quad t_i = 2i-1 \quad (1 \le i \le p),
\]
\[
  u_j = j \ (1 \le j < p),\qquad u_p = 0,
\]
and let the positive weights be
\begin{gather*}
  w_{R_iH_i} = r_i+h_i =: g_i,\qquad w_{R_iT_i} = r_i,\qquad w_{R_iU_j} = r_i+u_j,\\
  w_{S_iH_j} = s_i+h_j,\qquad w_{S_iT_i} = t_i,\qquad w_{S_1T_2} = s_1,
\end{gather*}
for all $i, j$; every other weight is $0$. Let $O_p$ be defined in the same way with
$r_i = 2p+i+2$, $s_i = 2i+2$, $t_i = 2i+1$ (and $h, u$ unchanged), plus a row $S_0$ and a column
$T_0$ with $w_{S_0H_j} = h_j+2$ for all $j$, $w_{S_0T_0} = 1$ and $w_{S_0T_1} = 2$ (so
$m = 2p+1$ and $|B| = 3p+1$). Then:
\begin{enumerate}
\item[(a)] the optimal face is a single point $f$, with $f = r$ on the $R_i$, $s$ on the $S_i$,
  $h$ on the $H_j$, $0$ on the $T_j$ and $u$ on the $U_j$, and for $O_p$ also $f_{S_0} = 2$,
  $f_{T_0} = 0$;
\item[(b)] every weakly-descending order fires exactly the pairs $(R_i,H_i)$ and $(S_i,T_i)$, and
  for $O_p$ also $(S_0,T_0)$, so the cold output $q$ is the same under every order;
\item[(c)] $\|q-y^*\|_0 = 5p-1 = \lfloor \tfrac52 m \rfloor - 1$ on $E_p$, and
  $\|q-y^*\|_0 = 5p+1 = \lfloor \tfrac52 m \rfloor - 1$ on $O_p$.
\end{enumerate}
\end{proposition}

\begin{proof}
(a) \emph{Optimality.} Every positive pair is covered by $f$, with equality at $(R_i,H_i)$,
$(R_i,T_i)$, all $(R_i,U_j)$, all $(S_i,H_j)$ and $(S_1,T_2)$ (and, for $O_p$, all $(S_0,H_j)$
and $(S_0,T_1)$). The remaining positive pairs are covered strictly:
$f_{S_i} = s_i = t_i+1$ against $w_{S_iT_i} = t_i$, and $f_{S_0} = 2 > 1 = w_{S_0T_0}$. The
matching $M^* = \{(R_i,U_i),(S_i,H_i)\}$, together with $(S_0,T_1)$ for $O_p$, is tight, so
$\sum f = w(M^*)$ and $f$ is optimal by weak duality.

\emph{Uniqueness.} Let $f'$ be optimal and $\delta = f'-f$. By complementary slackness $f'$ is
tight on $M^*$ and vanishes on every vertex that $M^*$ leaves unsaturated. For $E_p$ these are
$T_1,\dots,T_p$; for $O_p$ they are $T_0,T_2,\dots,T_p$. Every $(R_i,U_j)$ is tight at $f$, so
$\delta_{R_i}+\delta_{U_j} \ge 0$ for all $i, j$, with equality when $j = i$. Hence
$\delta_{R_i} \ge -\delta_{U_j} = \delta_{R_j}$ for all $i, j$, so $\delta_R \equiv \alpha$ and
$\delta_U \equiv -\alpha$. In the same way, the pairs $(S_i,H_j)$ with $i, j \ge 1$ give
$\delta_S \equiv \beta$ and $\delta_H \equiv -\beta$. Since $f_{U_p} = 0$, $\alpha \le 0$.

For $E_p$: $(R_1,T_1)$ is tight and $\delta_{T_1} = 0$, so $\alpha \ge 0$. $(S_1,T_2)$ is tight
and $\delta_{T_2} = 0$, so $\beta \ge 0$. $(R_1,H_1)$ is tight, so $\alpha-\beta \ge 0$. Hence
$\alpha = \beta = 0$ and $\delta = 0$.

For $O_p$: let $\tau = \delta_{T_1} \ge 0$. Tightness on $(S_0,T_1) \in M^*$ gives
$\delta_{S_0} = -\tau$, and the tight pairs $(S_0,H_j)$ give $-\tau-\beta \ge 0$. The pair
$(S_1,T_2)$ gives $\beta \ge 0$, so $\beta = \tau = 0$ and $\delta_{S_0} = 0$. Finally
$(R_2,T_2)$ is tight with $\delta_{T_2} = 0$, so $\alpha \ge 0$, hence $\alpha = 0$ and
$\delta = 0$. (Both arguments use $T_2$, hence $p \ge 2$.)

(b) \emph{Dominance.} Call the listed pairs \emph{intended}. They are positive and form a matching
that covers every row. Every positive pair $e$ that is not intended has an endpoint whose intended
pair is strictly heavier than $e$:
\begin{itemize}
\item $(R_i,T_i)$ and $(R_i,U_j)$: $w_e \le r_i+p-1 < r_i+h_i = g_i$ (endpoint $R_i$);
\item $(S_i,H_j)$: $w_e = s_i+h_j < r_j+h_j = g_j$ (endpoint $H_j$), because $s_i \le 2p < 2p+j = r_j$
  on $E_p$ and $s_i \le 2p+2 < 2p+j+2 = r_j$ on $O_p$ (this is why $O_p$ shifts $r$ by $2$);
\item $(S_1,T_2)$: $s_1 < t_2$, that is $2 < 3$ on $E_p$ and $4 < 5$ on $O_p$ (endpoint $T_2$);
\item on $O_p$, $(S_0,H_j)$: $h_j+2 < g_j$ (endpoint $H_j$), and $(S_0,T_1)$: $2 < 3 = t_1$
  (endpoint $T_1$).
\end{itemize}
Such a condition fixes the cold output. By induction along any weakly-descending order, when a
pair is processed every vertex whose intended pair came earlier carries that pair's weight, and
every other vertex is at $0$. An intended pair therefore meets two vertices at $0$ and fires. A
pair that is not intended has an endpoint whose intended pair is strictly heavier, so that pair
was processed earlier; the endpoint already exceeds $w_e$ and the pair does not fire. Zero pairs
never fire. So $q = g_i$ on $R_i$ and $H_i$, $q = t_i$ on $S_i$ and $T_i$, $q = 1$ on $S_0$ and
$T_0$ (for $O_p$), and $q = 0$ on every $U_j$.

(c) By (a), $\|q-y^*\|_0 = \|q-f\|_0$. The two disagree at $R_i$ ($g_i \ne r_i$), $H_i$
($g_i \ne h_i$), $S_i$ ($t_i = s_i - 1$) and $T_i$ ($t_i > 0$) for every $i$, and at $U_j$ for
$j < p$ ($u_j > 0$); they agree at $U_p$. That is $4p + (p-1) = 5p-1$ on $E_p$. On $O_p$, add
$S_0$ ($1 \ne 2$) and $T_0$ ($1 \ne 0$), for $5p+1$. Finally,
$\lfloor \tfrac52 \cdot 2p \rfloor - 1 = 5p-1$ and
$\lfloor \tfrac52 (2p+1) \rfloor - 1 = 5p+1$.
\end{proof}

\begin{remark}[The constant $\tfrac52$]\label{rem:coldconstant}
Theorem~\ref{thm:coldfivehalves}(iii) gives
$E(m)/m = \tfrac52 - \tfrac1m$ for even $m \ge 4$ and $\tfrac52 - \tfrac{3}{2m}$ for odd
$m \ge 3$. Hence $\sup_m E(m)/m = \tfrac52$ is approached but never attained.

The witnesses differ in how they depend on the tie-break. The $6\times8$ instance of
Proposition~\ref{prop:coldrefuted}(c) reaches face distance $14$ under the lexicographic
tie-break, but not under every order: its six cold outputs over all weakly-descending orders lie at
face distances $11, 11, 12, 12, 13, 14$, so its minimum over orders is $11$. The families $E_p$ and $O_p$, the
$3\times4$ witness and the $2\times3$ block of part (i) have a single cold output under every
weakly-descending order. For example, $E_3$ is a $6\times9$ instance at face distance $14$ under
every order, with a unique optimal dual.
\end{remark}

\subsection{Warm start}

\begin{lemma}[Firing chain]\label{lem:chain}
Let $|A| = 1$ and let $\repair_\downarrow$ fire at $k_1, k_2, \dots, k_m$ in processing order. Then
the \emph{initial} values (after the clamping $\hat{y} \leftarrow \max(\hat{y},0)$) satisfy
$\hat{y}_{k_1} > \hat{y}_{k_2} > \dots > \hat{y}_{k_m} \ge 0$.
\end{lemma}
\begin{proof}
Each coordinate is processed once, and the scan only raises, so when $k_i$ is processed its value is
still $\hat{y}_{k_i}$. A firing there sets
$y_a \leftarrow y_a + (w_{k_i} - y_a - \hat{y}_{k_i}) = w_{k_i} - \hat{y}_{k_i}$.
Since $y_a$ strictly increases at each firing,
$w_{k_i} - \hat{y}_{k_i} < w_{k_{i+1}} - \hat{y}_{k_{i+1}}$, and descending order gives
$w_{k_i} \ge w_{k_{i+1}}$, so
$\hat{y}_{k_{i+1}} < \hat{y}_{k_i} + (w_{k_{i+1}} - w_{k_i}) \le \hat{y}_{k_i}$.
\end{proof}

\begin{corollary}\label{cor:onenew}
For $|A| = 1$, $\repair_\downarrow$ creates \emph{at most one} new nonzero among the right-side
coordinates. (The left coordinate $y_a$ may also become nonzero; the count $N$ in the proof of
Proposition~\ref{prop:sortedwarm} does not include it.) The hypothesis is needed: on
$|A| = |B| = 2$ with $w_{00} = 5$, $w_{11} = 3$, $w_{01} = w_{10} = 1$ and $\hat y = 0$, the
descending scan fires twice and creates two.
\end{corollary}
\begin{proof}
By Lemma~\ref{lem:chain}, the values $\hat{y}_{k_1}, \dots, \hat{y}_{k_{m-1}}$ are all strictly
greater than $\hat{y}_{k_m} \ge 0$, hence all nonzero. So at most one fired index, the last,
can have started at $0$, and a firing at an already-nonzero coordinate creates no new nonzero.
\end{proof}

\begin{proof}[Proof of Proposition~\ref{prop:sortedwarm}]
Write $W = \max_k w_k$, let $W_2$ be the second largest weight counted with multiplicity (read
$W_2 = 0$ when $|B| = 1$, where there is no second weight), and fix an
argmax $k^*$. For $u = t \ge 0$ the best completion has value
$\mathrm{val}(t) = t + \sum_k (w_k - t)^+$, which equals $W$ exactly for $t \in [W_2, W]$ and
exceeds $W$ elsewhere; so $\opt = W$ and the optimal face is
$\{(t \mid v(t)) : t \in [W_2, W]\}$ with $v(t)_{k^*} = W - t$ and $v(t)_k = 0$ for $k \ne k^*$. Since
the coordinates off $k^*$ contribute a term independent of $t$, the distance decomposes as
$d(y) = N(y) + \delta(y)$, where $N(y) = |\{k \ne k^* : y_k \ne 0\}|$ and
$\delta(y) = \min_{t \in [W_2, W]} \big([y_a \ne t] + [y_{k^*} \ne W - t]\big)$.
Corollary~\ref{cor:onenew} gives $\Delta N \le 1$, so it suffices to prove $\Delta \delta \le 0$.
Note first that after any firing at $k$ we have $y_a = w_k - \hat{y}_k \le w_k \le W$, and if no
firing occurs then $y$, hence $\delta$, is unchanged; so we may assume some constraint fires, and
then the final iterate satisfies $y_a \le W$, which is all that is used below.

\emph{Case $\delta(\hat{y}) = 0$.} Then $\hat{y}_a \in [W_2, W]$ and $\hat{y}_{k^*} = W - \hat{y}_a$.
For $k \ne k^*$, $\hat{y}_a + \hat{y}_k \ge \hat{y}_a \ge W_2 \ge w_k$; and
$\hat{y}_a + \hat{y}_{k^*} = W = w_{k^*}$. So $\hat{y}$ is already feasible, no constraint fires, and
$\delta$ is unchanged.

\emph{Case $\delta(\hat{y}) = 2$.} Immediate, since $\delta \le 2$ always.

\emph{Case $\delta(\hat{y}) = 1$.} Suppose $\delta(\repair_\downarrow(\hat{y})) = 2$. As $y_{k^*} \ge 0$,
this means $y_a \notin [W_2, W]$ and $y_{k^*} > W - W_2$; with $y_a \le W$ we get $y_a < W_2$. Note that
$y_{k^*}$ is still at its initial value when constraint $k^*$ is processed (it is touched only
there), and $y_a \ge \hat{y}_a$ at every moment. If $k^*$ fires, then immediately afterwards
$y_a = W - \hat{y}_{k^*}$, and $y_a$ only increases, so $W - \hat{y}_{k^*} < W_2$, i.e.
$\hat{y}_{k^*} > W - W_2$. Then $\delta(\hat{y}) = 1$ forces $\hat{y}_a \in [W_2, W]$, but firing
at $k^*$ requires $y_a + \hat{y}_{k^*} < W$ at that moment, hence
$\hat{y}_a + \hat{y}_{k^*} < W$, while $\hat{y}_a \ge W_2$ and $\hat{y}_{k^*} > W - W_2$ give
$\hat{y}_a + \hat{y}_{k^*} > W$, a contradiction. If instead $k^*$
does not fire, then $y_{k^*} = \hat{y}_{k^*}$ (each coordinate is processed once), so
$\hat{y}_{k^*} > W - W_2$; and $y_a \ge \hat{y}_a$ gives $\hat{y}_a < W_2$. Both coordinates are then
misaligned, i.e. $\delta(\hat{y}) = 2$, contradicting the case assumption.

Hence $\Delta\delta \le 0$ and $\Delta d = \Delta N + \Delta \delta \le 1$.
\end{proof}

The argument covers ties: if the maximum weight is attained more than once, then $W_2 = W$, the
face collapses to the single point $(W \mid 0)$, and every step of the proof goes through with
$[W_2, W] = \{W\}$. Only weak descent is used, in Lemma~\ref{lem:chain} and nowhere else, so the
bound holds for every tie-break. An exhaustive sweep covers $150{,}827{,}580$ instances
($|B| \le 5$, all weights in $\{0,\dots,4\}$, all predictions in $\{0,\dots,5\}^{|B|+1}$) under
every weakly-descending order, $7.19\times10^{8}$ (instance, order) pairs in total. It finds zero
violations, and the $+1$ is attained on $9{,}896{,}766$ of these pairs.

\paragraph{The $+1$ needs $\min(|A|,|B|) = 1$.} No such bound holds once $\min(|A|,|B|) \ge 2$: on
the $2\times3$ instance $w = (1\,2\,5;\,4\,5\,6)$ with $\hat y = (0,0 \mid 0,0,3)$, both weakly-descending
orders end at face distance $4$ from distance $2$, and on a $2\times6$ instance the increase reaches $3$.
Theorem~\ref{thm:sharpwarm}, proved next, gives the worst case.

\begin{proof}[Proof of Theorem~\ref{thm:sharpwarm}]
Write $p = \max(\hat{y},0)$ and $q = \repair_{\downarrow}(p)$, fix any optimal dual
$f = (u^* \mid v^*)$, and fix a maximum-weight matching $M$ (an integral optimal primal exists by
total unimodularity). Throughout, $(u \mid v)$ denotes the scan's current iterate and
$\mathrm{ham}(\cdot\,,\cdot)$ counts the coordinates where two vectors differ. Call a coordinate
\emph{aligned} if $p$ agrees with $f$ there, and
\emph{damaged} if it is aligned and $q$ differs from $f$ there. Coordinates never decrease and
change only by hosting a firing, so a damaged coordinate is exactly an aligned coordinate that
hosts a firing, and any coordinate's first firing occurs while it still holds its starting value.
Writing $D(f)$ for the number of damaged coordinates and $G(f) \ge 0$ for the number of
coordinates brought into agreement with $f$,
\[
\mathrm{ham}(q, f) \;=\; \mathrm{ham}(p, f) + D(f) - G(f) \;\le\; \mathrm{ham}(p, f) + D(f).
\]
Since $f \ge 0$, clamping helps coordinatewise, so it suffices to prove
$D(f) \le 2\min(|A|,|B|) - 1$ for every optimal dual $f$ and apply it to a face point attaining
$\|y^* - p\|_0$.

\emph{(a) Zero-current firings are terminal.} If $(a,b)$ fires while $v_b = 0$,
then afterwards $u_a = w_{ab}$, and every later pair $(a, b')$ has $w_{ab'} \le w_{ab} \le u_a$
with $u_a$ non-decreasing, so it cannot fire. Hence each row hosts at most one such firing, and it
is that row's last. \emph{(b) Complementary slackness.} If $v^*_b > 0$ then $b$ is matched in $M$,
say to row $r(b)$, with $u^*_{r(b)} + v^*_b = w_{r(b),b}$; symmetrically for rows.
\emph{(c) First firers are low.} If an aligned column $b$ is first fired via $(a', b)$, then
$v_b = v^*_b$ there, and $w_{a'b} - u_{a'} - v^*_b > 0$ with feasibility
$w_{a'b} \le u^*_{a'} + v^*_b$ force $u_{a'} < u^*_{a'}$ at that moment.
\emph{(d) Aligned rows never fire $f$-zero columns.} If $p_a = u^*_a$ then $u_a \ge u^*_a$
throughout, and any column with $v^*_b = 0$ has $w_{ab} \le u^*_a \le u_a$, residual $\le 0$.

\emph{Charging: $D(f) \le 2|A|$.} Give each row $a$ two slots. $S_1(a)$ is filled if $a$ is
damaged, or if $a$ is the first firer of a damaged column with $v^*_b = 0$; by (d) these cases
exclude each other, and by (a) distinct such columns have distinct first firers. $S_2(a)$ is
filled if $a$'s $M$-partner is an aligned, damaged column with $v^*_b > 0$; by (b) every such
column is matched, and since $M$ is a matching this is injective. Every damaged coordinate fills
a distinct slot, so $D(f) \le 2|A|$.

\emph{Refinement: not all $2|A|$ slots can fill.} Suppose they all do. Then every row is either
damaged (aligned, $u_a \ge u^*_a$ throughout, never low: an \emph{H-row}) or the first firer of a
damaged $f$-zero column (an \emph{L-row}), exclusively; and every row's $M$-partner $c_a$ is an
aligned, positive, damaged column. An L-row $\ell$ fires its zero column $z_\ell$ last, and there
$u_\ell < w_{\ell z_\ell} \le u^*_\ell$ (feasibility at $v^*_{z_\ell} = 0$), so $\ell$ is low at
all times up to that firing. The set $L$ of L-rows is nonempty: any filled $S_2$ names a damaged
aligned positive column, whose first firer is low by (c), hence not an H-row. Let $g(a)$ be the
first firer of $c_a$; by (c), $g$ maps into $L$.

$g(\ell) \ne \ell$ for $\ell \in L$: if $\ell$ first fired $c_\ell$ itself, tightness (b) makes
that residual $u^*_\ell - u_\ell$, so afterwards $u_\ell = u^*_\ell$; the pair $(\ell, z_\ell)$
has weight $w_{\ell z_\ell} \le u^*_\ell < u^*_\ell + v^*_{c_\ell} = w_{\ell c_\ell}$, so it is
processed strictly later, and its residual there is $\le 0$, contradicting that $\ell$ zero-fires
(the zero-firing cannot come first, since it is terminal). So $g$ is fixed-point-free on the nonempty
finite set $L$; iterating yields a cycle
$\rho_1 \mapsto \rho_2 \mapsto \dots \mapsto \rho_m \mapsto \rho_1$ with $m \ge 2$. Consider an
edge $\rho \mapsto \tau = g(\rho)$. The pair $(\rho, c_\rho)$ has weight
$w_{\rho c_\rho} = u^*_\rho + v^*_{c_\rho} > u^*_\rho \ge w_{\rho z_\rho}$, so it strictly
precedes $\rho$'s zero pair, and $\rho$ is still low there; were $c_\rho$ then untouched, $\rho$
would fire it and be its own first firer. Hence $\tau$'s first firing of $c_\rho$ strictly
precedes the pair $(\rho, c_\rho)$, which in a weakly-descending scan forces
$w_{\tau c_\rho} \ge w_{\rho c_\rho}$; with feasibility
$w_{\tau c_\rho} \le u^*_\tau + v^*_{c_\rho}$ and tightness
$w_{\rho c_\rho} = u^*_\rho + v^*_{c_\rho}$ this gives $u^*_\tau \ge u^*_\rho$. Around the cycle
all these hold with equality, so $w_{\tau c_\rho} = u^*_\tau + v^*_{c_\rho}$: $\tau$'s first
firing of $c_\rho$ has residual $u^*_\tau - u_\tau$ and sets $u_\tau = u^*_\tau$, after which
$\tau$'s strictly later zero pair cannot fire, contradicting $\tau \in L$.
So the slots cannot all fill, and $D(f) \le 2|A| - 1$.

\emph{Transpose.} The operator and every step above are symmetric in the two sides, and a
weakly-descending order remains weakly descending under transposition, so also
$D(f) \le 2|B| - 1$. The argument used
only that values never decrease, each pair is processed once, weak descent, feasibility, and
complementary slackness; neither integrality nor the tie-break enters.
\end{proof}

\subsection{The balanced family}

Throughout, $\mathcal{B}_n$, the dual $f = (u^* \mid v^*)$ and the prediction $\hat y = f - n\,e_{a_1}$
are as in Definition~\ref{def:balancedfam}.

\begin{lemma}[The optimal face of $\mathcal{B}_n$]\label{lem:balancedface}
$\opt(\mathcal{B}_n) = n^2$, and the optimal face is the triangle
\[
  F \;=\; \bigl\{\, f(m,\gamma) \;:\; m \ge 0,\ \gamma \ge 0,\ m + \gamma \le 1 \,\bigr\},
\]
\[
  f(m,\gamma) \;=\; \bigl(n - m,\ \gamma, \dots, \gamma \;\big|\; 2(n-1) - \gamma,\ \dots,\ 2 - \gamma,\ m\bigr),
\]
that is, $f(m,\gamma)_{a_1} = n - m$, $f(m,\gamma)_{a_i} = \gamma$ for $i \ge 2$, $f(m,\gamma)_{b_k} = 2(n-k) - \gamma$ for $k < n$, and $f(m,\gamma)_{b_n} = m$. Its vertices are $f = f(0,0)$, $f(1,0)$ and $f(0,1)$, and these are its only integral points.
\end{lemma}
\begin{proof}
$f \ge 0$ is feasible: for $i \ge 2$, $u^*_{a_i} + v^*_{b_k} = 2(n-k) = w_{a_i b_k}$, and
$u^*_{a_1} + v^*_{b_k} = 3n-2k \ge 2n-k$ because $k \le n$. Its value is $n + \sum_k 2(n-k) = n^2$,
and the matching $M = \{(a_1,b_n)\} \cup \{(a_i,b_{i-1})\}_{i \ge 2}$ has weight $n + n(n-1) = n^2$,
so $\opt = n^2$ and $f$ is optimal.

Now let $y = (u \mid v)$ be any optimal dual and write $m = n - u_{a_1}$, $\gamma_i = u_{a_i}$ for
$i \ge 2$, $\Gamma = \sum_{i \ge 2} \gamma_i$, and $\beta_k = v_{b_k} - 2(n-k)$. Optimality and
feasibility give the following.
\begin{itemize}
\item[(S)] $-m + \Gamma + \sum_k \beta_k = 0$, since $y$ and $f$ both have value $\opt$.
\item[(P)] $u \ge 0$ and $v \ge 0$; its row half is (N) $\gamma_i \ge 0$.
\item[(T$_i$)] $\gamma_i + \beta_k \ge 0$ for $i \ge 2$ and every $k$: feasibility where $f$ is
  tight.
\item[(T$_1$)] $\beta_k \ge k - n + m$: feasibility at $(a_1, b_k)$.
\end{itemize}
By (T$_1$) at $k = n$, $\beta_n \ge m$, so (S) gives $\sum_{k<n} \beta_k = m - \Gamma - \beta_n \le
-\Gamma$. By (T$_i$) at a row $i$ of minimum $\gamma_i$, $\beta_k \ge -\min_i \gamma_i$ for every
$k$, so $\sum_{k<n} \beta_k \ge -(n-1)\min_i \gamma_i \ge -\Gamma$. Hence equality holds throughout:
the $\gamma_i$ share a common value $\gamma$, $\beta_k = -\gamma$ for every $k < n$, and
$\beta_n = m$.

Three constraints remain: $\gamma \ge 0$ from (N), $m \ge 0$ from $v_{b_n} \ge 0$ in (P), and
$m + \gamma \le 1$ from (T$_1$) at $k = n-1$. The other $k < n$ in (T$_1$) are
weaker, and the rest of (P) and of (T$_i$) follow from these three and $n \ge 2$. Conversely, every
$f(m,\gamma)$ with $m, \gamma \ge 0$ and $m + \gamma \le 1$ is feasible. After the substitution,
(T$_i$) reads $0 \ge 0$ and $\gamma + m \ge 0$, (T$_1$) reads $m + \gamma \le n - k$, and (P)
follows from $m + \gamma \le 1$. Each such $f(m,\gamma)$ has value
$n - m + (n-1)\gamma + \sum_{k<n}\bigl(2(n-k) - \gamma\bigr) + m = n^2$. The coordinates of
$f(m,\gamma)$ are integral exactly when $m$ and $\gamma$ are, which on the triangle leaves its
three vertices.
\end{proof}

\begin{proof}[Proof of Theorem~\ref{thm:balancedlb}]
\emph{Step 1: $\opt = n^2$ and $f$ is optimal.} This is Lemma~\ref{lem:balancedface}.

\emph{Step 2: $\|y^*-\hat y\|_0 = 1$.} $\hat y$ differs from $f$ in one coordinate, and it is not
on the face, since it is infeasible: $\hat y_{a_1} + \hat y_{b_1} = 2n-2 < 2n-1 = w_{a_1 b_1}$.

\emph{Step 3: every weakly-descending order gives the same output.} Coordinates never decrease, so
$y_{b_k} \ge 2(n-k)$ throughout; hence for $i \ge 2$ the residual at $(a_i,b_k)$ is at most
$2(n-k) - 0 - 2(n-k) = 0$ and rows $a_2,\dots,a_n$ never fire. So $y_{a_1}$ moves only in row $a_1$
and $y_{b_k}$ only at $(a_1,b_k)$. The weights $w_{a_1 b_k} = 2n-k$ are strictly decreasing
in $k$, so every weakly-descending order processes $(a_1,b_1),\dots,(a_1,b_n)$ in that relative
order; no tie-break can reorder them. By induction, when $(a_1,b_k)$ is reached $y_{a_1} = k-1$
and $y_{b_k} = 2(n-k)$, so the residual is $(2n-k)-(k-1)-(2n-2k) = 1 > 0$: it fires, leaving
$y_{a_1} = k$ and $y_{b_k} = 2(n-k)+1$. The output is $q$ with $q_{a_1} = n$, $q_{a_i} = 0$ and
$q_{b_k} = 2(n-k)+1$.

\emph{Step 4: $\|y^*-q\|_0 = n$.} By Lemma~\ref{lem:balancedface} every face point is some
$f(m,\gamma)$ with $m, \gamma \ge 0$, $m + \gamma \le 1$, and
\[
  q - f(m,\gamma) \;=\; \bigl(m,\ -\gamma, \dots, -\gamma \;\big|\; 1+\gamma,\ \dots,\ 1+\gamma,\ 1-m\bigr).
\]
The $n-1$ columns $b_k$ with $k < n$ always differ, as $1 + \gamma \ge 1$; the rows $a_i$ with
$i \ge 2$ differ exactly when $\gamma \ne 0$; and $a_1$ and $b_n$ differ exactly when $m \ne 0$ and
$m \ne 1$ respectively, at least one of which holds. So $\mathrm{ham}(q, f(m,\gamma)) \ge (n-1) + 1
= n$, with equality at $\gamma = 0$ and $m \in \{0, 1\}$.
\end{proof}

\begin{proof}[Proof of Corollary~\ref{cor:noadditive}]
By Theorem~\ref{thm:balancedlb}, for every $n \ge 2$ the instance $\mathcal{B}_n$ has
$\min(|A|,|B|) = n$, a prediction at face distance $1$, and an output at face distance $n$ under
every weakly-descending order. That is an increase of $n - 1 = \min(|A|,|B|) - 1$, so no
additive term below $\min(|A|,|B|) - 1$ is valid, and the optimal constant on balanced $n \times n$
instances is at least $n - 1$. The upper end $2n - 1$ is the additive term of
Theorem~\ref{thm:sharpwarm} at $\min(|A|,|B|) = n$.
\end{proof}

\begin{remark}[A balanced companion to Theorem~\ref{thm:order}]\label{rem:balancedorder}
On $\mathcal{B}_n$ the weakly-increasing scan fires exactly once, at $(a_1,b_n)$ with residual
$n$, and ends at $q' = f + n\,e_{b_n}$, at face distance $1$. So in $\ell_0$ the two orders end at
$n$ versus $1$ while agreeing on the gap and both $\ell_1$ distances, as in Theorem~\ref{thm:order}: both
outputs have dual gap $n$, realised $\ell_1$ displacement $2n$, and $\ell_1$ distance $n$ to the
optimal face. The $\ell_\infty$ reversal of Section~\ref{sec:delicate} (Appendix~\ref{app:order})
recurs here: the descending output lies at $\ell_\infty$ distance $1$ from the face, the increasing
output at $n-1$. By Lemma~\ref{lem:balancedface}, $q - f(m,\gamma)$ has sup norm $1 + \gamma$,
minimised at $\gamma = 0$, while
$q' - f(m,\gamma) = (m, -\gamma, \dots, -\gamma \mid \gamma, \dots, \gamma, n-m)$ has sup norm
$n - m \ge n - 1$, attained at $m = 1$. The same two minima hold in $\ell_\infty^\pm$ in covering
coordinates, where the values are $1 + 2\gamma$ and $n - m + \gamma$; in the difference coordinates
of \citet{sakaue2022} the increasing output is at $n$ and the descending one at $1$, so the reversal
persists. For $n \ge 3$ the $\ell_\infty$ ranking is therefore strict and again reverses the
$\ell_0$ one. Theorem~\ref{thm:order} has $|A| = 1$ and the forward scan loses, while $\mathcal{B}_n$ is balanced
and the sorted scan loses. Neither order is always better, so $\mathcal{B}_n$ is not a lower bound
over all orders. No order is additive on covering instances (Theorem~\ref{thm:noadditiveclass}), but
a data-dependent order, the cover-directed scan, stays within twice the prediction's $\ell_0$ error
(Theorem~\ref{thm:coverscan}).
\end{remark}

\paragraph{Transfer to the packing orientation.} \citet[Thm.~3.2]{chen2022} work in the packing
orientation, and not everything transfers to it. Under $z_i = T - y_i$ and $c_{ab} = 2T - w_{ab}$ a
covering system becomes a packing one. The map is a coordinatewise bijection, hence an $\ell_0$
isometry, and it carries residuals to excesses. Non-negative predictions, such as that of
$\mathcal{B}_n$, leave the clamp of Definition~\ref{def:greedy} inert. On them $\repair_\sigma$ is
therefore carried exactly onto the packing mirror, and a weakly-descending weight order becomes a
weakly-ascending capacity one. $\mathcal{B}_n$ is balanced with a perfect optimal matching, so its
image is a minimum-cost perfect-matching instance, and the balanced lower bound of
Theorem~\ref{thm:balancedlb} carries over to the packing orientation.

The perfect-matching dual has no non-negativity constraint, and dropping it changes neither $\opt$
nor the two face distances. The perfect matching of Lemma~\ref{lem:balancedface} still certifies
$\opt = n^2$. That lemma's proof reaches $\gamma_i = \gamma$, $\beta_k = -\gamma$ for $k < n$ and
$\beta_n = m$ without (P), so the face becomes the unbounded strip
$\{f(m,\gamma) : m, \gamma \in \R,\ 0 \le m + \gamma \le 1\}$, which still contains $f$. In Step~4
of the proof of Theorem~\ref{thm:balancedlb}, the rows $a_i$, $i \ge 2$, agree with $q$ only at
$\gamma = 0$ and the columns $b_k$, $k < n$, only at $\gamma = -1$. So at least $n - 1$ of them
differ, and $a_1$, $b_n$ add at least one more. Equality holds at $f$, which is still on the face,
and $\hat y$ stays at distance $1$.

The family of Theorem~\ref{thm:order} does not transfer: with $|A| = 1$ and $|B| = d \ge 2$ it has
no perfect matching.

\subsection{Sharpness in the unbalanced regime}

In the unbalanced regime the additive constant of Theorem~\ref{thm:sharpwarm} cannot be lowered for
any value of $\min(|A|,|B|)$, as the following family shows.
\begin{definition}[The extremal family $\mathcal{F}_k$]\label{def:sharpfam}
Fix $k \ge 1$, put $M = 4k+2$ and $s = 2k+1$, and for $1 \le i \le k-1$ put
\[
  p_i = M - 2i - 1,\qquad q_i = M - 2i - 2,\qquad h_i = 2i+1,\qquad \ell_i = 2i .
\]
$\mathcal{F}_k$ has $A = \{0,1,\dots,k-1\}$ (row $0$ the \emph{hub}, the rest \emph{spokes}) and
$3k+2$ columns
\[
  X_1,\; X_2,\; L_1, \dots, L_{k-1},\; Z_0,\; H_1, Z_1, \dots, H_{k-1}, Z_{k-1},\; P_1,\; P_2 .
\]
All weights are $0$ except
\begin{gather*}
  w_{0X_1} = M,\quad w_{0X_2} = M-1,\quad w_{0Z_0} = s,\\
  w_{0L_i} = p_i,\quad w_{iL_i} = q_i,\quad w_{iH_i} = h_i,\quad w_{iZ_i} = \ell_i ,
\end{gather*}
so the padding columns $P_1, P_2$ carry no positive weight. The prediction is
$\hat y_0 = 0$, $\hat y_i = \ell_i - 1$, $\hat y_{X_1} = M$, $\hat y_{X_2} = M-1$,
$\hat y_{L_i} = q_i - h_i$, $\hat y_{H_i} = 2$, and $0$ on every $Z$ and $P$. So
$\min(|A|,|B|) = k$, $\bigl||B|-|A|\bigr| = 2k+2$ and $\opt = 3k^2 + k + 2$.
\end{definition}

\begin{proposition}[The constant is attained at every $\min(|A|,|B|)$]\label{prop:sharpattained}
For every $k \ge 1$, on $\mathcal{F}_k$ and under \emph{every} weakly-descending scan order,
\[
  \|y^* - \hat y\|_0 = 2k+1, \qquad \|y^* - \repair_{\downarrow}(\hat y)\|_0 = 4k ,
\]
both minima over the optimal face; hence
$\|y^* - \repair_{\downarrow}(\hat y)\|_0 = \|y^* - \hat y\|_0 + 2\min(|A|,|B|) - 1$. Also,
$2k+1 < 2k+2 = \bigl||B|-|A|\bigr|$ and $4k < 4k+2 = |A|+|B|$: the family lies where the bound
has content, and the additive term of the minimum binds.
\end{proposition}

\begin{proof}
Write $g_0 = 0$ and $g_i = h_i + 1 = 2i+2$.

\emph{The order is unique.} The positive weights are $M$, $M-1$, the odd $p_i \in \{2k+3,\dots,4k-1\}$,
the even $q_i \in \{2k+2,\dots,4k-2\}$, the odd $s = 2k+1$, the odd $h_i \in \{3,\dots,2k-1\}$ and the
even $\ell_i \in \{2,\dots,2k-2\}$: six pairwise disjoint sets, so all positive weights are
distinct. A pair of weight $0$ has residual $-y_a - y_b \le 0$ throughout and never fires. So every
weakly-descending order induces the same firings and $\repair_\downarrow(\hat y)$ is tie-break
independent.

\emph{The optimal face.} Feasibility on $\{0,X_1\}$ and on each $\{i, L_i\}$ gives
$\sum_V y \ge M + \sum_i q_i$, and the matching $\{(0,X_1)\} \cup \{(i,L_i)\}$ attains it using the
largest entry of every row, so $\opt = M + \sum_i q_i = 3k^2+k+2$. Equality forces
$y_{X_2} = y_{Z_0} = y_{H_i} = y_{Z_i} = y_{P_j} = 0$ and both groups tight; then $y_0 \ge M-1$ and
$y_i \ge h_i$, while $y_0 + y_{L_i} \ge p_i$ holds automatically as
$y_i \le q_i \le 4k-2 < 4k+1 \le y_0$. So
\[
  F = \bigl\{\, y_0 \in [M-1, M],\; y_{X_1} = M - y_0,\;
      y_i \in [h_i, q_i],\; y_{L_i} = q_i - y_i,\; y \equiv 0 \text{ elsewhere} \,\bigr\},
\]
the full, possibly fractional, face.

\emph{Before the repair.} Every $y \in F$ disagrees with $\hat y$ at $y_0$ ($y_0 \ge M-1 > 0$), at
$y_{X_1}$ ($\le 1 < M$), at $y_{X_2}$ ($0 \ne M-1$), at each $y_i$ ($y_i \ge h_i > \ell_i - 1$) and
at each $y_{H_i}$ ($0 \ne 2$): $2k+1$ forced disagreements. The point with $y_0 = M$, $y_i = h_i$
agrees on the other $2k+1$ coordinates, so $\|y^* - \hat y\|_0 = 2k+1$.

\emph{The scan fires exactly $2k-1$ times.} $(0,X_1)$ and $(0,X_2)$ come first and are tight
at $\hat y$, so they never fire, and $y_0$ only increases. Since
$p_1 > \dots > p_{k-1} > s$, the pairs $(0,L_i)$ are processed in increasing $i$ before $(0,Z_0)$;
inductively $(0,L_i)$ meets $y_0 = g_{i-1}$, $y_{L_i} = q_i - h_i$ and fires with residual
$g_i - g_{i-1} > 0$. The pair $(i,L_i)$, of smaller weight $q_i < p_i$, is then slack. Next
$(0,Z_0)$ is reached with $y_0 = g_{k-1}$ (that is, $2k$ for $k \ge 2$, and $g_0 = 0$ for
$k = 1$) and fires with residual $s - g_{k-1} \ge 1$. On a spoke, $(i,H_i)$ precedes $(i,Z_i)$ and
is tight at $\hat y$, so only $(i,Z_i)$ fires, with residual $1$.
Every other pair has weight $0$. So the firings are exactly
$(0,L_1),\dots,(0,L_{k-1}),(0,Z_0),(1,Z_1),\dots,(k-1,Z_{k-1})$: $2k-1$ of them.

\emph{After the repair.} The $2k+1$ coordinates already off the face stay off:
$X_1, X_2, H_i$ are untouched and $y_0$ ends at $s = 2k+1 < M-1$, $y_i$ at $\ell_i < h_i$. Each of
the $2k-1$ damaged coordinates is now off too: $y_{Z_0} = s - g_{k-1} > 0$, $y_{Z_i} = 1 \ne 0$ and
$y_{L_i} > q_i - h_i \ge q_i - y_i$ for every $y \in F$. Only $P_1, P_2$ still agree, so
$\|y^* - \repair_\downarrow(\hat y)\|_0 = (4k+2) - 2 = 4k$.
\end{proof}

We also verify $\mathcal{F}_k$ by computer (Appendix~\ref{app:repro}): through
$\min(|A|,|B|) = 12$ by exact minimisation of the Hamming distance over the full, possibly
fractional optimal face, and independently through $\min(|A|,|B|) = 5$ by brute-force enumeration
of the integral face. The two agree at every $k$ where both run, as the total unimodularity
argument of Appendix~\ref{app:early} requires.

\begin{remark}[Large prediction errors]
The warm-start bounds are additive in the prediction's $\ell_0$ error, so they help only when that
error is small. A $\hat{y}$ drawn uniformly at
random is far from the optimal face in $\ell_0$ to begin with; $\repair_\downarrow$ inherits that
distance, and the post-repair face distance grows with the instance. None of these bounds promises
$\ell_0$ sparsity that the prediction lacks.
\end{remark}
